%============================================================
%  Bd. 15 - Halbgruppen, Monoide und Halbringe
%============================================================

\documentclass[main.tex]{subfiles}

\ifSubfilesClassLoaded{
  \BandSetupStandalone{B15}
}{
}

\title{Bd. 15 - Halbgruppen, Monoide und Halbringe}
\author{Martin Kunze}
\date{}
\setcounter{file}{15}

\begin{document}

\title{Bd. 15 - Halbgruppen, Monoide und Halbringe}
\author{Martin Kunze}
\date{}
\hypersetup{pageanchor=false}
\maketitle
\hypersetup{pageanchor=true}
\setcounter{tocdepth}{3}
\tableofcontents
\setcounter{file}{15}
\renewcommand{\FormulaBandID}{15}

\chapter{Einführung}

Dieser Band abstrahiert die in Band~12 bewiesenen Rechengesetze der
natürlichen Zahlen zu Halbgruppen, Monoiden und Halbringen. Algebraische
Strukturen werden stets zusammen mit ihren Operationen und ausgezeichneten
Elementen bezeichnet.

\chapter{Halbgruppen}

\section{Definition der Halbgruppe}

\FormulaDefDeltaK[Begriff der Halbgruppe]{%
  \SemigroupAlg{A}{\star}%
}{SemigroupDef}{%
  \DeltaRow{Mengen}{A\dsep x\dsep y\dsep z}
  \DeltaRow{\textbf{Axiome}}{}
  \DeltaRow{Abgeschlossenheit}{%
    x,y\in A\vdash x\star y\in A}
  \DeltaRow{Assoziativität}{%
    x,y,z\in A\vdash (x\star y)\star z=x\star(y\star z)}
  \DeltaRow{\textbf{Neues Symbol}}{}
  \DeltaPrem{Halbgruppe}{\SemigroupAlg{A}{\star}}
}

\FormulaAxiomDeltaK[Zugriff auf die Abgeschlossenheit einer Halbgruppe]{%
  \SemigroupAlg{A}{\star}\dsep x,y\in A
  \vdash x\star y\in A%
}{SemigroupClosureAxiom}{%
  \DeltaRow{Mengen}{A\dsep x\dsep y}
}

\FormulaAxiomDeltaK[Zugriff auf die Assoziativität einer Halbgruppe]{%
  \SemigroupAlg{A}{\star}\dsep x,y,z\in A
  \vdash (x\star y)\star z=x\star(y\star z)%
}{SemigroupAssociativityAxiom}{%
  \DeltaRow{Mengen}{A\dsep x\dsep y\dsep z}
}

\FormulaAxiomDeltaK[Zugriff auf die binäre Operation einer Halbgruppe]{%
  \SemigroupAlg{A}{\star}\vdash
  \mathsf{BinOp}\!\left(A,\star\right)%
}{SemigroupBinaryOperationAxiom}{%
  \DeltaRow{Mengen}{A}
}

\section{Kommutative Halbgruppen}

\FormulaDefDeltaK[Begriff der kommutativen Halbgruppe]{%
  \CommutativeSemigroupAlg{A}{\star}%
}{CommutativeSemigroupDef}{%
  \DeltaRow{Mengen}{A\dsep x\dsep y}
  \DeltaRow{\textbf{Axiome}}{}
  \DeltaPrem{Halbgruppe}{\SemigroupAlg{A}{\star}}
  \DeltaRow{Kommutativität}{%
    x,y\in A\vdash x\star y=y\star x}
  \DeltaRow{\textbf{Neues Symbol}}{}
  \DeltaPrem{Kommutative Halbgruppe}{%
    \CommutativeSemigroupAlg{A}{\star}}
}

\FormulaAxiomDeltaK[Zugriff auf die Halbgruppenstruktur]{%
  \CommutativeSemigroupAlg{A}{\star}
  \vdash\SemigroupAlg{A}{\star}%
}{CommutativeSemigroupBaseAxiom}{%
  \DeltaRow{Mengen}{A}
}

\FormulaAxiomDeltaK[Zugriff auf die Kommutativität]{%
  \CommutativeSemigroupAlg{A}{\star}\dsep x,y\in A
  \vdash x\star y=y\star x%
}{CommutativeSemigroupCommutativityAxiom}{%
  \DeltaRow{Mengen}{A\dsep x\dsep y}
}

\FormulaThmDeltaK[Dreigliedriges Vertauschungsgesetz]{%
  \CommutativeSemigroupAlg{A}{\star}\dsep a,b,c\in A
  \vdash (a\star b)\star c=(a\star c)\star b%
}{CommutativeSemigroupThreeTermExchange}{%
  \DeltaRow{Mengen}{A\dsep a\dsep b\dsep c}
}
\begin{tabproofwide}
  \proofstepwidestar[1]{\CommutativeSemigroupAlg{A}{\star}}{\rA}
  \proofstepwidestar[2]{a\in A}{\rA}
  \proofstepwidestar[3]{b\in A}{\rA}
  \proofstepwidestar[4]{c\in A}{\rA}
  \proofstepwidestar[1]{\SemigroupAlg{A}{\star}}{%
    \FormulaRefAuto{CommutativeSemigroupBaseAxiom}[ax]{1}}
  \proofstepwidestar[1,2,3,4]{%
    (a\star c)\star b=a\star(c\star b)}{%
    \FormulaRefAuto{SemigroupAssociativityAxiom}[ax]{5,2,4,3}}
  \proofstepwide[1,2,3,4]{(a\star b)\star c}{=}{a\star(b\star c)}{%
    \FormulaRefAuto{SemigroupAssociativityAxiom}[ax]{5,2,3,4}}
  \proofstepwide[1,3,4]{}{=}{a\star(c\star b)}{%
    \FormulaRefAuto{CommutativeSemigroupCommutativityAxiom}[ax]{1,3,4}}
  \proofstepwide[1,2,3,4]{}{=}{(a\star c)\star b}{%
    \FormulaRefAuto{a=b\vdash b=a}{6}}
  \proofstepwidestar[1,2,3,4]{%
    (a\star b)\star c=(a\star c)\star b}{%
    \rChain{7,9}}
\end{tabproofwide}

\FormulaThmDeltaK[Viergliedriges Vertauschungsgesetz]{%
  \begin{aligned}[t]
    &\CommutativeSemigroupAlg{A}{\star}\dsep a,b,c,d\in A\vdash{}\\[-2pt]
    &(a\star b)\star(c\star d)=(a\star c)\star(b\star d)
  \end{aligned}
}{CommutativeSemigroupFourTermExchange}{%
  \DeltaRow{Mengen}{A\dsep a\dsep b\dsep c\dsep d}
}
\begin{tabproofwide}
  \proofstepwidestar[1]{\CommutativeSemigroupAlg{A}{\star}}{\rA}
  \proofstepwidestar[2]{a\in A}{\rA}
  \proofstepwidestar[3]{b\in A}{\rA}
  \proofstepwidestar[4]{c\in A}{\rA}
  \proofstepwidestar[5]{d\in A}{\rA}
  \proofstepwidestar[1]{\SemigroupAlg{A}{\star}}{%
    \FormulaRefAuto{CommutativeSemigroupBaseAxiom}[ax]{1}}
  \proofstepwidestar[1,2,3]{a\star b\in A}{%
    \FormulaRefAuto{SemigroupClosureAxiom}[ax]{6,2,3}}
  \proofstepwidestar[1,2,4]{a\star c\in A}{%
    \FormulaRefAuto{SemigroupClosureAxiom}[ax]{6,2,4}}
  \proofstepwidestar[1,2,3,4,5]{%
    \begin{aligned}[t]
      &(a\star b)\star(c\star d)\mathrel{=}\\[-2pt]
      &\qquad\bigl((a\star b)\star c\bigr)\star d
    \end{aligned}}{%
    \FormulaRefAuto{a=b\vdash b=a}{%
      \FormulaRefAuto{SemigroupAssociativityAxiom}[ax]{6,7,4,5}}}
  \proofstepwidestar[1,2,3,4,5]{%
    \begin{aligned}[t]
      &\bigl((a\star b)\star c\bigr)\star d\mathrel{=}\\[-2pt]
      &\qquad\bigl((a\star c)\star b\bigr)\star d
    \end{aligned}}{%
    \FormulaRefAuto{CommutativeSemigroupThreeTermExchange}{1,2,3,4}}
  \proofstepwidestar[1,2,3,4,5]{%
    \begin{aligned}[t]
      &\bigl((a\star c)\star b\bigr)\star d\mathrel{=}\\[-2pt]
      &\qquad(a\star c)\star(b\star d)
    \end{aligned}}{%
    \FormulaRefAuto{SemigroupAssociativityAxiom}[ax]{6,8,3,5}}
  \proofstepwidestar[1,2,3,4,5]{%
    (a\star b)\star(c\star d)=(a\star c)\star(b\star d)}{%
    \rChain{9,11}}
\end{tabproofwide}

\section{Idempotente Halbgruppen}

\FormulaDefDeltaK[Begriff der idempotenten Halbgruppe]{%
  \mathsf{IdemHGr}\!\left(A,\star\right)%
}{IdempotentSemigroupDef}{%
  \DeltaRow{Mengen}{A\dsep x}
  \DeltaRow{\textbf{Axiome}}{}
  \DeltaPrem{Halbgruppe}{\SemigroupAlg{A}{\star}}
  \DeltaRow{Idempotenz}{x\in A\vdash x\star x=x}
  \DeltaRow{\textbf{Neues Symbol}}{}
  \DeltaPrem{Idempotente Halbgruppe}{%
    \mathsf{IdemHGr}\!\left(A,\star\right)}
}

\FormulaAxiomDeltaK[Zugriff auf die Halbgruppe einer idempotenten Halbgruppe]{%
  \mathsf{IdemHGr}\!\left(A,\star\right)
  \vdash\SemigroupAlg{A}{\star}%
}{IdempotentSemigroupBaseAxiom}{%
  \DeltaRow{Mengen}{A}
}

\FormulaAxiomDeltaK[Zugriff auf die Idempotenz]{%
  \mathsf{IdemHGr}\!\left(A,\star\right)\dsep x\in A
  \vdash x\star x=x%
}{IdempotentSemigroupIdempotenceAxiom}{%
  \DeltaRow{Mengen}{A\dsep x}
}

\FormulaDefDeltaK[Kommutative idempotente Halbgruppe]{%
  \mathsf{KIHGr}\!\left(A,\star\right)%
}{CommutativeIdempotentSemigroupDef}{%
  \DeltaRow{Mengen}{A\dsep x}
  \DeltaRow{\textbf{Axiome}}{}
  \DeltaPrem{Kommutative Halbgruppe}{%
    \CommutativeSemigroupAlg{A}{\star}}
  \DeltaRow{Idempotenz}{x\in A\vdash x\star x=x}
  \DeltaRow{\textbf{Neues Symbol}}{}
  \DeltaPrem{\makecell[l]{Kommutative idempotente\\ Halbgruppe}}{%
    \mathsf{KIHGr}\!\left(A,\star\right)}
}

\FormulaAxiomDeltaK[Zugriff auf die kommutative Halbgruppe]{%
  \mathsf{KIHGr}\!\left(A,\star\right)
  \vdash\CommutativeSemigroupAlg{A}{\star}%
}{CommutativeIdempotentSemigroupBaseAxiom}{%
  \DeltaRow{Mengen}{A}
}

\FormulaAxiomDeltaK[Zugriff auf die Idempotenz der kommutativen Struktur]{%
  \mathsf{KIHGr}\!\left(A,\star\right)\dsep x\in A
  \vdash x\star x=x%
}{CommutativeIdempotentSemigroupIdempotenceAxiom}{%
  \DeltaRow{Mengen}{A\dsep x}
}

\section{Kürzbare Operationen}

\FormulaDefDeltaK[Linkskürzbare Operation]{%
  \mathsf{LK\ddot{u}rz}\!\left(A,\star\right)%
}{LeftCancellativeDef}{%
  \DeltaRow{Mengen}{A\dsep a\dsep b\dsep c}
  \DeltaRow{Linkskürzung}{%
    a,b,c\in A\dsep a\star b=a\star c\vdash b=c}
  \DeltaRow{\textbf{Neues Symbol}}{}
  \DeltaPrem{Linkskürzbarkeit}{\mathsf{LK\ddot{u}rz}\!\left(A,\star\right)}
}

\FormulaDefDeltaK[Rechtskürzbare Operation]{%
  \mathsf{RK\ddot{u}rz}\!\left(A,\star\right)%
}{RightCancellativeDef}{%
  \DeltaRow{Mengen}{A\dsep a\dsep b\dsep c}
  \DeltaRow{Rechtskürzung}{%
    a,b,c\in A\dsep b\star a=c\star a\vdash b=c}
  \DeltaRow{\textbf{Neues Symbol}}{}
  \DeltaPrem{Rechtskürzbarkeit}{\mathsf{RK\ddot{u}rz}\!\left(A,\star\right)}
}

\FormulaDefDeltaK[Kürzbare Operation]{%
  \mathsf{K\ddot{u}rz}\!\left(A,\star\right)%
}{CancellativeOperationDef}{%
  \DeltaRow{Mengen}{A}
  \DeltaPrem{Linkskürzbarkeit}{\mathsf{LK\ddot{u}rz}\!\left(A,\star\right)}
  \DeltaPrem{Rechtskürzbarkeit}{\mathsf{RK\ddot{u}rz}\!\left(A,\star\right)}
  \DeltaRow{\textbf{Neues Symbol}}{}
  \DeltaPrem{Kürzbarkeit}{\mathsf{K\ddot{u}rz}\!\left(A,\star\right)}
}

\FormulaAxiomDeltaK[Zugriff auf die Linkskürzung]{%
  \mathsf{LK\ddot{u}rz}\!\left(A,\star\right)\dsep
  a,b,c\in A\dsep a\star b=a\star c
  \vdash b=c%
}{LeftCancellationAxiom}{%
  \DeltaRow{Mengen}{A\dsep a\dsep b\dsep c}
}

\FormulaAxiomDeltaK[Zugriff auf die Rechtskürzung]{%
  \mathsf{RK\ddot{u}rz}\!\left(A,\star\right)\dsep
  a,b,c\in A\dsep b\star a=c\star a
  \vdash b=c%
}{RightCancellationAxiom}{%
  \DeltaRow{Mengen}{A\dsep a\dsep b\dsep c}
}

\FormulaAxiomDeltaK[Eine kürzbare Operation ist linkskürzbar]{%
  \mathsf{K\ddot{u}rz}\!\left(A,\star\right)
  \vdash\mathsf{LK\ddot{u}rz}\!\left(A,\star\right)%
}{CancellativeLeftAxiom}{%
  \DeltaRow{Mengen}{A}
}

\FormulaAxiomDeltaK[Eine kürzbare Operation ist rechtskürzbar]{%
  \mathsf{K\ddot{u}rz}\!\left(A,\star\right)
  \vdash\mathsf{RK\ddot{u}rz}\!\left(A,\star\right)%
}{CancellativeRightAxiom}{%
  \DeltaRow{Mengen}{A}
}

\FormulaThmDeltaK[Rechtskürzungsinstanz aus Linkskürzung und Kommutativität]{%
  \begin{aligned}[t]
    &\CommutativeSemigroupAlg{A}{\star}\dsep
      \mathsf{LK\ddot{u}rz}\!\left(A,\star\right)\dsep
      a,b,c\in A\dsep b\star a=c\star a\vdash{}\\[-2pt]
    &b=c
  \end{aligned}
}{CommutativeLeftCancellationImpliesRightCancellation}{%
  \DeltaRow{Mengen}{A\dsep a\dsep b\dsep c}
}
\begin{tabproofwide}
  \proofstepwidestar[1]{\CommutativeSemigroupAlg{A}{\star}}{\rA}
  \proofstepwidestar[2]{\mathsf{LK\ddot{u}rz}\!\left(A,\star\right)}{\rA}
  \proofstepwidestar[3]{a\in A}{\rA}
  \proofstepwidestar[4]{b\in A}{\rA}
  \proofstepwidestar[5]{c\in A}{\rA}

  \proofstepwide[1,3,4]{a\star b}{=}{b\star a}{%
    \FormulaRefAuto{CommutativeSemigroupCommutativityAxiom}[ax]{1,3,4}}
  \proofstepwide[7]{}{=}{c\star a}{\rA}
  \proofstepwide[1,3,5]{}{=}{a\star c}{%
    \FormulaRefAuto{CommutativeSemigroupCommutativityAxiom}[ax]{1,5,3}}
  \proofstepwidestar[1,3,4,5,7]{a\star b=a\star c}{%
    \rChain{6,8}}

  \proofstepwidestar[1,2,3,4,5,7]{b=c}{%
    \FormulaRefAuto{LeftCancellationAxiom}[ax]{2,3,4,5,9}}
\end{tabproofwide}

\FormulaThmDeltaK[Kommutative linkskürzbare Halbgruppen sind rechtskürzbar]{%
  \CommutativeSemigroupAlg{A}{\star}\dsep
  \mathsf{LK\ddot{u}rz}\!\left(A,\star\right)
  \vdash\mathsf{RK\ddot{u}rz}\!\left(A,\star\right)%
}{CommutativeLeftCancellativeImpliesRightCancellative}{%
  \DeltaRow{Mengen}{A\dsep a\dsep b\dsep c}
}
\begin{tabproof}
  \proofstep{1}{\CommutativeSemigroupAlg{A}{\star}}{\rA}
  \proofstep{2}{\mathsf{LK\ddot{u}rz}\!\left(A,\star\right)}{\rA}
  \proofstep{3}{a\in A}{\rA}
  \proofstep{4}{b\in A}{\rA}
  \proofstep{5}{c\in A}{\rA}
  \proofstep{6}{b\star a=c\star a}{\rA}
  \proofstep{1,2,3,4,5,6}{b=c}{%
    \FormulaRefAuto{CommutativeLeftCancellationImpliesRightCancellation}[thm]{%
      1,2,3,4,5,6}}
  \proofstep{1,2}{\mathsf{RK\ddot{u}rz}\!\left(A,\star\right)}{%
    \FormulaRefAuto{RightCancellativeDef}[def]{7}}
\end{tabproof}

\section{Morphismen von Halbgruppen}

\FormulaDefDeltaK[Begriff des Halbgruppenhomomorphismus]{%
  \SemigroupHom{f}{A,\star}{B,\diamond}%
}{SemigroupHomDef}{%
  \DeltaRow{Mengen}{A\dsep B\dsep x\dsep y}
  \DeltaRow{Funktionen}{f\dsep\star\dsep\diamond}
  \DeltaRow{\textbf{Axiome}}{}
  \DeltaPrem{Quellhalbgruppe}{\SemigroupAlg{A}{\star}}
  \DeltaPrem{Zielhalbgruppe}{\SemigroupAlg{B}{\diamond}}
  \DeltaPrem{Abbildung}{f\colon A\to B}
  \DeltaRow{Operationserhaltung}{%
    x,y\in A\vdash
    f(x\star y)=f(x)\diamond f(y)}
  \DeltaRow{\textbf{Neues Symbol}}{}
  \DeltaPrem{Halbgruppenhomomorphismus}{%
    \SemigroupHom{f}{A,\star}{B,\diamond}}
}

\FormulaAxiomDeltaK[Operationserhaltung eines Halbgruppenhomomorphismus]{%
  \SemigroupHom{f}{A,\star}{B,\diamond}\dsep x,y\in A
  \vdash f(x\star y)=f(x)\diamond f(y)%
}{SemigroupHomOperationAxiom}{%
  \DeltaRow{Mengen}{A\dsep B\dsep x\dsep y}
  \DeltaRow{Funktionen}{f\dsep\star\dsep\diamond}
}

\FormulaAxiomDeltaK[Trägerabbildung eines Halbgruppenhomomorphismus]{%
  \SemigroupHom{f}{A,\star}{B,\diamond}
  \vdash f\colon A\to B%
}{SemigroupHomFunctionAxiom}{%
  \DeltaRow{Mengen}{A\dsep B}
  \DeltaRow{Funktionen}{f\dsep\star\dsep\diamond}
}

\FormulaDefDeltaK[Begriff des Halbgruppenisomorphismus]{%
  \SemigroupIso{f}{A,\star}{B,\diamond}%
}{SemigroupIsoDef}{%
  \DeltaRow{Mengen}{A\dsep B}
  \DeltaRow{Funktionen}{f\dsep\star\dsep\diamond}
  \DeltaRow{\textbf{Axiome}}{}
  \DeltaPrem{Halbgruppenhomomorphismus}{%
    \SemigroupHom{f}{A,\star}{B,\diamond}}
  \DeltaPrem{Bijektivität}{f\colon A\bij B}
  \DeltaRow{\textbf{Neues Symbol}}{}
  \DeltaPrem{Halbgruppenisomorphismus}{%
    \SemigroupIso{f}{A,\star}{B,\diamond}}
}

\subsection{Grundlegende Eigenschaften}

\FormulaThmDeltaK[Komposition von Halbgruppenhomomorphismen]{%
  \begin{aligned}[t]
    &\SemigroupHom{f}{A,\star}{B,\diamond}\dsep
      \SemigroupHom{g}{B,\diamond}{C,\bullet}\\
    &\vdash\SemigroupHom{g\circ f}{A,\star}{C,\bullet}
  \end{aligned}
}{SemigroupHomComposition}{%
  \DeltaRow{Mengen}{A\dsep B\dsep C\dsep x\dsep y}
  \DeltaRow{Funktionen}{f\dsep g\dsep\star\dsep\diamond\dsep\bullet}
}
\begin{tabproofwide}
  \proofstepwidestar[1]{\SemigroupHom{f}{A,\star}{B,\diamond}}{\rA}
  \proofstepwidestar[2]{\SemigroupHom{g}{B,\diamond}{C,\bullet}}{\rA}
  \proofstepwidestar[1]{f\colon A\to B}{%
    \FormulaRefAuto{SemigroupHomFunctionAxiom}[ax]{1}}
  \proofstepwidestar[2]{g\colon B\to C}{%
    \FormulaRefAuto{SemigroupHomFunctionAxiom}[ax]{2}}
  \proofstepwidestar[1,2]{g\circ f\colon A\to C}{%
    \FormulaRefAuto{G\circ F\colon A\to C}[thm]{3,4}}
  \proofstepwidestar[1]{\SemigroupAlg{A}{\star}}{%
    \FormulaRefAuto{SemigroupHomDef}[def]{1}}
  \proofstepwidestar[2]{\SemigroupAlg{C}{\bullet}}{%
    \FormulaRefAuto{SemigroupHomDef}[def]{2}}
  \proofstepwidestar[8]{x\in A}{\rA}
  \proofstepwidestar[9]{y\in A}{\rA}
  \proofstepwidestar[1,8,9]{x\star y\in A}{%
    \FormulaRefAuto{SemigroupClosureAxiom}[ax]{6,8,9}}
  \proofstepwidestar[1,8]{f(x)\in B}{%
    \FormulaRefAuto{%
      F\colon A\to B\dsep x\in A\vdash F(x)\in B%
    }[thm]{3,8}}
  \proofstepwidestar[1,9]{f(y)\in B}{%
    \FormulaRefAuto{%
      F\colon A\to B\dsep x\in A\vdash F(x)\in B%
    }[thm]{3,9}}
  \proofstepwidestar[1,8,9]{%
    f(x\star y)=f(x)\diamond f(y)}{%
    \FormulaRefAuto{SemigroupHomOperationAxiom}[ax]{1,8,9}}
  \proofstepwidestar[1,2,8]{g(f(x))=(g\circ f)(x)}{%
    \FormulaRefAuto{%
      x\in A \vdash F(G(x))=(F\circ G)(x)%
    }[thm]{3,4,8}}
  \proofstepwidestar[1,2,9]{g(f(y))=(g\circ f)(y)}{%
    \FormulaRefAuto{%
      x\in A \vdash F(G(x))=(F\circ G)(x)%
    }[thm]{3,4,9}}
  \proofstepwide[1,2,8,9]{(g\circ f)(x\star y)}{=}{%
    g(f(x\star y))}{%
    \FormulaRefAuto{%
      x\in A \vdash (F\circ G)(x)=F(G(x))%
    }[thm]{3,4,10}}
  \proofstepwide[1,2,8,9]{}{=}{g(f(x)\diamond f(y))}{%
    \rIE{13}}
  \proofstepwide[1,2,8,9]{}{=}{g(f(x))\bullet g(f(y))}{%
    \FormulaRefAuto{SemigroupHomOperationAxiom}[ax]{2,11,12}}
  \proofstepwide[1,2,8,9]{}{=}{%
    (g\circ f)(x)\bullet g(f(y))}{\rIE{14}}
  \proofstepwide[1,2,8,9]{}{=}{%
    (g\circ f)(x)\bullet(g\circ f)(y)}{\rIE{15}}
  \proofstepwide[1,2,8,9]{(g\circ f)(x\star y)}{=}{%
    (g\circ f)(x)\bullet(g\circ f)(y)}{%
    \rChain{16,20}}
  \proofstepwidestar[1,2]{%
    \SemigroupHom{g\circ f}{A,\star}{C,\bullet}}{%
    \FormulaRefAuto{SemigroupHomDef}[def]{6,7,5,21}}
\end{tabproofwide}

\FormulaThmDeltaK[Die Identität ist ein Halbgruppenisomorphismus]{%
  \SemigroupAlg{A}{\star}
  \vdash\SemigroupIso{\Id_A}{A,\star}{A,\star}%
}{SemigroupIdentityIsomorphism}{%
  \DeltaRow{Mengen}{A\dsep x\dsep y}
  \DeltaRow{Funktionen}{\star}
}
\begin{tabproofwide}
  \proofstepwidestar[1]{\SemigroupAlg{A}{\star}}{\rA}
  \proofstepwidestar[]{\Id_A\colon A\to A}{%
    \FormulaRefAuto{IdA}[thm]}
  \proofstepwidestar[2]{x\in A}{\rA}
  \proofstepwidestar[3]{y\in A}{\rA}
  \proofstepwidestar[1,2,3]{x\star y\in A}{%
    \FormulaRefAuto{SemigroupClosureAxiom}[ax]{1,2,3}}
  \proofstepwidestar[2]{\Id_A(x)=x}{%
    \FormulaRefAuto{x\in A\vdash \Id_A(x)=x}{2}}
  \proofstepwidestar[2]{x=\Id_A(x)}{%
    \FormulaRefAuto{a=b\vdash b=a}{6}}
  \proofstepwidestar[3]{\Id_A(y)=y}{%
    \FormulaRefAuto{x\in A\vdash \Id_A(x)=x}{3}}
  \proofstepwidestar[3]{y=\Id_A(y)}{%
    \FormulaRefAuto{a=b\vdash b=a}{8}}
  \proofstepwide[1,2,3]{\Id_A(x\star y)}{=}{x\star y}{%
    \FormulaRefAuto{x\in A\vdash \Id_A(x)=x}{5}}
  \proofstepwide[1,2,3]{}{=}{\Id_A(x)\star y}{\rIE{7}}
  \proofstepwide[1,2,3]{}{=}{\Id_A(x)\star\Id_A(y)}{\rIE{9}}
  \proofstepwidestar[1,2,3]{%
    \Id_A(x\star y)=\Id_A(x)\star\Id_A(y)}{\rChain{10,12}}
  \proofstepwidestar[1]{%
    \SemigroupHom{\Id_A}{A,\star}{A,\star}}{%
    \FormulaRefAuto{SemigroupHomDef}[def]{1,1,2,13}}
  \proofstepwidestar[]{\Id_A\colon A\bij A}{%
    \FormulaRefAuto{\Id_A\colon A\bij A}[thm]}
  \proofstepwidestar[1]{%
    \SemigroupIso{\Id_A}{A,\star}{A,\star}}{%
    \FormulaRefAuto{SemigroupIsoDef}[def]{14,15}}
\end{tabproofwide}

\FormulaThmDeltaK[Komposition von Halbgruppenisomorphismen]{%
  \begin{aligned}[t]
    &\SemigroupIso{f}{A,\star}{B,\diamond}\dsep
      \SemigroupIso{g}{B,\diamond}{C,\bullet}\\
    &\vdash\SemigroupIso{g\circ f}{A,\star}{C,\bullet}
  \end{aligned}
}{SemigroupIsoComposition}{%
  \DeltaRow{Mengen}{A\dsep B\dsep C}
  \DeltaRow{Funktionen}{f\dsep g\dsep\star\dsep\diamond\dsep\bullet}
}
\begin{tabproof}
  \proofstep{1}{\SemigroupIso{f}{A,\star}{B,\diamond}}{\rA}
  \proofstep{2}{\SemigroupIso{g}{B,\diamond}{C,\bullet}}{\rA}
  \proofstep{1}{\SemigroupHom{f}{A,\star}{B,\diamond}}{%
    \FormulaRefAuto{SemigroupIsoDef}[def]{1}}
  \proofstep{2}{\SemigroupHom{g}{B,\diamond}{C,\bullet}}{%
    \FormulaRefAuto{SemigroupIsoDef}[def]{2}}
  \proofstep{1,2}{\SemigroupHom{g\circ f}{A,\star}{C,\bullet}}{%
    \FormulaRefAuto{SemigroupHomComposition}[thm]{3,4}}
  \proofstep{1}{f\colon A\bij B}{%
    \FormulaRefAuto{SemigroupIsoDef}[def]{1}}
  \proofstep{2}{g\colon B\bij C}{%
    \FormulaRefAuto{SemigroupIsoDef}[def]{2}}
  \proofstep{1,2}{g\circ f\colon A\bij C}{%
    \FormulaRefAuto{F\colon A\bij B \dsep G\colon B\bij C\vdash
      G\circ F\colon A\bij C}[thm]{6,7}}
  \proofstep{1,2}{\SemigroupIso{g\circ f}{A,\star}{C,\bullet}}{%
    \FormulaRefAuto{SemigroupIsoDef}[def]{5,8}}
\end{tabproof}

\subsection{Spezialisierte Homomorphismen und Isomorphismen}

\FormulaDefDeltaK[Spezialisierte Halbgruppenhomomorphismen]{%
  \begin{aligned}[t]
    &\CommutativeSemigroupHom{f}{A,\star}{B,\diamond}\\[-2pt]
    &\quad\coloneqq
      \CommutativeSemigroupAlg{A}{\star}\land
      \CommutativeSemigroupAlg{B}{\diamond}\\[-2pt]
    &\qquad\land\SemigroupHom{f}{A,\star}{B,\diamond},\\[2pt]
    &\IdempotentSemigroupHom{f}{A,\star}{B,\diamond}\\[-2pt]
    &\quad\coloneqq
      \mathsf{IdemHGr}(A,\star)\land
      \mathsf{IdemHGr}(B,\diamond)\\[-2pt]
    &\qquad\land\SemigroupHom{f}{A,\star}{B,\diamond},\\[2pt]
    &\CommutativeIdempotentSemigroupHom{f}{A,\star}{B,\diamond}\\[-2pt]
    &\quad\coloneqq
      \mathsf{KIHGr}(A,\star)\land
      \mathsf{KIHGr}(B,\diamond)\\[-2pt]
    &\qquad\land\SemigroupHom{f}{A,\star}{B,\diamond},\\[2pt]
    &\CancellativeSemigroupHom{f}{A,\star}{B,\diamond}\\[-2pt]
    &\quad\coloneqq
      \mathsf{K\ddot{u}rz}(A,\star)\land
      \mathsf{K\ddot{u}rz}(B,\diamond)\\[-2pt]
    &\qquad\land\SemigroupHom{f}{A,\star}{B,\diamond}.
  \end{aligned}
}{SpecializedSemigroupHomDef}{%
  \DeltaRow{Mengen}{A\dsep B}
  \DeltaRow{Funktionen}{f\dsep\star\dsep\diamond}
  \DeltaRow{Erhaltungsgesetz}{%
    x,y\in A\vdash f(x\star y)=f(x)\diamond f(y)}
  \DeltaRow{\textbf{Neue Symbole}}{}
  \DeltaPrem{Kommutativer Fall}{\mathsf{KHGrHom}}
  \DeltaPrem{Idempotenter Fall}{\mathsf{IdemHGrHom}}
  \DeltaPrem{Kommutativ-idempotenter Fall}{\mathsf{KIHGrHom}}
  \DeltaPrem{Kürzbarer Fall}{\mathsf{K\ddot{u}rzHGrHom}}
}

\FormulaAxiomDeltaK[Quellstruktur eines Homomorphismus kommutativer Halbgruppen]{%
  \CommutativeSemigroupHom{f}{A,\star}{B,\diamond}
  \vdash\CommutativeSemigroupAlg{A}{\star}%
}{CommutativeSemigroupHomSourceAxiom}{%
  \DeltaRow{Mengen}{A\dsep B}
  \DeltaRow{Funktionen}{f\dsep\star\dsep\diamond}
}

\FormulaAxiomDeltaK[Zielstruktur eines Homomorphismus kommutativer Halbgruppen]{%
  \CommutativeSemigroupHom{f}{A,\star}{B,\diamond}
  \vdash\CommutativeSemigroupAlg{B}{\diamond}%
}{CommutativeSemigroupHomTargetAxiom}{%
  \DeltaRow{Mengen}{A\dsep B}
  \DeltaRow{Funktionen}{f\dsep\star\dsep\diamond}
}

\FormulaAxiomDeltaK[Halbgruppenhomomorphismus eines kommutativen Homomorphismus]{%
  \CommutativeSemigroupHom{f}{A,\star}{B,\diamond}
  \vdash\SemigroupHom{f}{A,\star}{B,\diamond}%
}{CommutativeSemigroupHomBaseHomAxiom}{%
  \DeltaRow{Mengen}{A\dsep B}
  \DeltaRow{Funktionen}{f\dsep\star\dsep\diamond}
}

\FormulaAxiomDeltaK[Quellstruktur eines Homomorphismus idempotenter Halbgruppen]{%
  \IdempotentSemigroupHom{f}{A,\star}{B,\diamond}
  \vdash\mathsf{IdemHGr}(A,\star)%
}{IdempotentSemigroupHomSourceAxiom}{%
  \DeltaRow{Mengen}{A\dsep B}
  \DeltaRow{Funktionen}{f\dsep\star\dsep\diamond}
}

\FormulaAxiomDeltaK[Zielstruktur eines Homomorphismus idempotenter Halbgruppen]{%
  \IdempotentSemigroupHom{f}{A,\star}{B,\diamond}
  \vdash\mathsf{IdemHGr}(B,\diamond)%
}{IdempotentSemigroupHomTargetAxiom}{%
  \DeltaRow{Mengen}{A\dsep B}
  \DeltaRow{Funktionen}{f\dsep\star\dsep\diamond}
}

\FormulaAxiomDeltaK[Halbgruppenhomomorphismus eines idempotenten Homomorphismus]{%
  \IdempotentSemigroupHom{f}{A,\star}{B,\diamond}
  \vdash\SemigroupHom{f}{A,\star}{B,\diamond}%
}{IdempotentSemigroupHomBaseHomAxiom}{%
  \DeltaRow{Mengen}{A\dsep B}
  \DeltaRow{Funktionen}{f\dsep\star\dsep\diamond}
}

\FormulaAxiomDeltaK[Quellstruktur eines kommutativ-idempotenten Homomorphismus]{%
  \CommutativeIdempotentSemigroupHom{f}{A,\star}{B,\diamond}
  \vdash\mathsf{KIHGr}(A,\star)%
}{CommutativeIdempotentSemigroupHomSourceAxiom}{%
  \DeltaRow{Mengen}{A\dsep B}
  \DeltaRow{Funktionen}{f\dsep\star\dsep\diamond}
}

\FormulaAxiomDeltaK[Zielstruktur eines kommutativ-idempotenten Homomorphismus]{%
  \CommutativeIdempotentSemigroupHom{f}{A,\star}{B,\diamond}
  \vdash\mathsf{KIHGr}(B,\diamond)%
}{CommutativeIdempotentSemigroupHomTargetAxiom}{%
  \DeltaRow{Mengen}{A\dsep B}
  \DeltaRow{Funktionen}{f\dsep\star\dsep\diamond}
}

\FormulaAxiomDeltaK[Halbgruppenhomomorphismus eines kommutativ-idempotenten Homomorphismus]{%
  \CommutativeIdempotentSemigroupHom{f}{A,\star}{B,\diamond}
  \vdash\SemigroupHom{f}{A,\star}{B,\diamond}%
}{CommutativeIdempotentSemigroupHomBaseHomAxiom}{%
  \DeltaRow{Mengen}{A\dsep B}
  \DeltaRow{Funktionen}{f\dsep\star\dsep\diamond}
}

\FormulaAxiomDeltaK[Quellstruktur eines Homomorphismus kürzbarer Halbgruppen]{%
  \CancellativeSemigroupHom{f}{A,\star}{B,\diamond}
  \vdash\mathsf{K\ddot{u}rz}(A,\star)%
}{CancellativeSemigroupHomSourceAxiom}{%
  \DeltaRow{Mengen}{A\dsep B}
  \DeltaRow{Funktionen}{f\dsep\star\dsep\diamond}
}

\FormulaAxiomDeltaK[Zielstruktur eines Homomorphismus kürzbarer Halbgruppen]{%
  \CancellativeSemigroupHom{f}{A,\star}{B,\diamond}
  \vdash\mathsf{K\ddot{u}rz}(B,\diamond)%
}{CancellativeSemigroupHomTargetAxiom}{%
  \DeltaRow{Mengen}{A\dsep B}
  \DeltaRow{Funktionen}{f\dsep\star\dsep\diamond}
}

\FormulaAxiomDeltaK[Halbgruppenhomomorphismus eines kürzbaren Homomorphismus]{%
  \CancellativeSemigroupHom{f}{A,\star}{B,\diamond}
  \vdash\SemigroupHom{f}{A,\star}{B,\diamond}%
}{CancellativeSemigroupHomBaseHomAxiom}{%
  \DeltaRow{Mengen}{A\dsep B}
  \DeltaRow{Funktionen}{f\dsep\star\dsep\diamond}
}

\FormulaDefDeltaK[Spezialisierte Halbgruppenisomorphismen]{%
  \begin{aligned}[t]
    &\CommutativeSemigroupIso{f}{A,\star}{B,\diamond}\\[-2pt]
    &\quad\coloneqq
      \CommutativeSemigroupHom{f}{A,\star}{B,\diamond}
      \land f\colon A\bij B,\\[2pt]
    &\IdempotentSemigroupIso{f}{A,\star}{B,\diamond}\\[-2pt]
    &\quad\coloneqq
      \IdempotentSemigroupHom{f}{A,\star}{B,\diamond}
      \land f\colon A\bij B,\\[2pt]
    &\CommutativeIdempotentSemigroupIso{f}{A,\star}{B,\diamond}\\[-2pt]
    &\quad\coloneqq
      \CommutativeIdempotentSemigroupHom{f}{A,\star}{B,\diamond}
      \land f\colon A\bij B,\\[2pt]
    &\CancellativeSemigroupIso{f}{A,\star}{B,\diamond}\\[-2pt]
    &\quad\coloneqq
      \CancellativeSemigroupHom{f}{A,\star}{B,\diamond}
      \land f\colon A\bij B.
  \end{aligned}
}{SpecializedSemigroupIsoDef}{%
  \DeltaRow{Mengen}{A\dsep B}
  \DeltaRow{Funktionen}{f\dsep\star\dsep\diamond}
  \DeltaRow{\textbf{Neue Symbole}}{}
  \DeltaPrem{Kommutativer Fall}{\mathsf{KHGrIso}}
  \DeltaPrem{Idempotenter Fall}{\mathsf{IdemHGrIso}}
  \DeltaPrem{Kommutativ-idempotenter Fall}{\mathsf{KIHGrIso}}
  \DeltaPrem{Kürzbarer Fall}{\mathsf{K\ddot{u}rzHGrIso}}
}

\FormulaAxiomDeltaK[Homomorphismus eines Isomorphismus kommutativer Halbgruppen]{%
  \CommutativeSemigroupIso{f}{A,\star}{B,\diamond}
  \vdash\CommutativeSemigroupHom{f}{A,\star}{B,\diamond}%
}{CommutativeSemigroupIsoHomAxiom}{%
  \DeltaRow{Mengen}{A\dsep B}
  \DeltaRow{Funktionen}{f\dsep\star\dsep\diamond}
}

\FormulaAxiomDeltaK[Bijektivität eines Isomorphismus kommutativer Halbgruppen]{%
  \CommutativeSemigroupIso{f}{A,\star}{B,\diamond}
  \vdash f\colon A\bij B%
}{CommutativeSemigroupIsoBijectiveAxiom}{%
  \DeltaRow{Mengen}{A\dsep B}
  \DeltaRow{Funktionen}{f\dsep\star\dsep\diamond}
}

\FormulaAxiomDeltaK[Homomorphismus eines Isomorphismus idempotenter Halbgruppen]{%
  \IdempotentSemigroupIso{f}{A,\star}{B,\diamond}
  \vdash\IdempotentSemigroupHom{f}{A,\star}{B,\diamond}%
}{IdempotentSemigroupIsoHomAxiom}{%
  \DeltaRow{Mengen}{A\dsep B}
  \DeltaRow{Funktionen}{f\dsep\star\dsep\diamond}
}

\FormulaAxiomDeltaK[Bijektivität eines Isomorphismus idempotenter Halbgruppen]{%
  \IdempotentSemigroupIso{f}{A,\star}{B,\diamond}
  \vdash f\colon A\bij B%
}{IdempotentSemigroupIsoBijectiveAxiom}{%
  \DeltaRow{Mengen}{A\dsep B}
  \DeltaRow{Funktionen}{f\dsep\star\dsep\diamond}
}

\FormulaAxiomDeltaK[Homomorphismuszugriff]{%
  \CommutativeIdempotentSemigroupIso{f}{A,\star}{B,\diamond}
  \vdash\CommutativeIdempotentSemigroupHom{f}{A,\star}{B,\diamond}%
}{CommutativeIdempotentSemigroupIsoHomAxiom}{%
  \DeltaRow{Mengen}{A\dsep B}
  \DeltaRow{Funktionen}{f\dsep\star\dsep\diamond}
}

\FormulaAxiomDeltaK[Bijektivität eines kommutativ-idempotenten Isomorphismus]{%
  \CommutativeIdempotentSemigroupIso{f}{A,\star}{B,\diamond}
  \vdash f\colon A\bij B%
}{CommutativeIdempotentSemigroupIsoBijectiveAxiom}{%
  \DeltaRow{Mengen}{A\dsep B}
  \DeltaRow{Funktionen}{f\dsep\star\dsep\diamond}
}

\FormulaAxiomDeltaK[Homomorphismus eines Isomorphismus kürzbarer Halbgruppen]{%
  \CancellativeSemigroupIso{f}{A,\star}{B,\diamond}
  \vdash\CancellativeSemigroupHom{f}{A,\star}{B,\diamond}%
}{CancellativeSemigroupIsoHomAxiom}{%
  \DeltaRow{Mengen}{A\dsep B}
  \DeltaRow{Funktionen}{f\dsep\star\dsep\diamond}
}

\FormulaAxiomDeltaK[Bijektivität eines Isomorphismus kürzbarer Halbgruppen]{%
  \CancellativeSemigroupIso{f}{A,\star}{B,\diamond}
  \vdash f\colon A\bij B%
}{CancellativeSemigroupIsoBijectiveAxiom}{%
  \DeltaRow{Mengen}{A\dsep B}
  \DeltaRow{Funktionen}{f\dsep\star\dsep\diamond}
}

\chapter{Monoide}

\section{Definition des Monoids}

\FormulaDefDeltaK[Begriff des Monoids]{%
  \MonoidAlg{A}{\star}{e}%
}{MonoidDef}{%
  \DeltaRow{Mengen}{A\dsep e\dsep x}
  \DeltaRow{\textbf{Axiome}}{}
  \DeltaPrem{Halbgruppe}{\SemigroupAlg{A}{\star}}
  \DeltaRow{Neutrales Element im Träger}{e\in A}
  \DeltaRow{Linksneutralität}{x\in A\vdash e\star x=x}
  \DeltaRow{Rechtsneutralität}{x\in A\vdash x\star e=x}
  \DeltaRow{\textbf{Neues Symbol}}{}
  \DeltaPrem{Monoid}{\MonoidAlg{A}{\star}{e}}
}

\FormulaAxiomDeltaK[Zugriff auf die Halbgruppe eines Monoids]{%
  \MonoidAlg{A}{\star}{e}\vdash\SemigroupAlg{A}{\star}%
}{MonoidBaseSemigroupAxiom}{%
  \DeltaRow{Mengen}{A\dsep e}
}

\FormulaAxiomDeltaK[Zugriff auf das neutrale Element]{%
  \MonoidAlg{A}{\star}{e}\vdash e\in A%
}{MonoidIdentityCarrierAxiom}{%
  \DeltaRow{Mengen}{A\dsep e}
}

\FormulaAxiomDeltaK[Zugriff auf die Linksneutralität]{%
  \MonoidAlg{A}{\star}{e}\dsep x\in A
  \vdash e\star x=x%
}{MonoidLeftIdentityAxiom}{%
  \DeltaRow{Mengen}{A\dsep e\dsep x}
}

\FormulaAxiomDeltaK[Zugriff auf die Rechtsneutralität]{%
  \MonoidAlg{A}{\star}{e}\dsep x\in A
  \vdash x\star e=x%
}{MonoidRightIdentityAxiom}{%
  \DeltaRow{Mengen}{A\dsep e\dsep x}
}

\section{Das neutrale Element}

\FormulaThmDeltaK[Ein neutraler Kandidat stimmt mit dem neutralen Element überein]{%
  \begin{aligned}[t]
    &\MonoidAlg{A}{\star}{e}\dsep e'\in A\dsep
      \forall x\in A\,(e'\star x=x)\dsep
      \forall x\in A\,(x\star e'=x)\vdash{}\\[-2pt]
    &e=e'
  \end{aligned}
}{MonoidIdentityCandidateEqualsIdentity}{%
  \DeltaRow{Mengen}{A\dsep e\dsep e'\dsep x}
}
\begin{tabproofwide}
    \proofstepwidestar[1]{\MonoidAlg{A}{\star}{e}}{\rA}
    \proofstepwidestar[2]{e'\in A}{\rA}
    \proofstepwidestar[3]{\forall x\in A\,(e'\star x=x)}{\rA}
    \proofstepwidestar[4]{\forall x\in A\,(x\star e'=x)}{\rA}
    \proofstepwidestar[1]{e\in A}{%
      \FormulaRefAuto{MonoidIdentityCarrierAxiom}[ax]{1}}
    \proofstepwidestar[1,3]{e'\star e=e}{\rRE{5,3}}
    \proofstepwide[1,3]{e}{=}{e'\star e}{%
      \FormulaRefAuto{a=b\vdash b=a}{6}}
    \proofstepwide[1,2]{}{=}{e'}{%
      \FormulaRefAuto{MonoidRightIdentityAxiom}[ax]{1,2}}
    \proofstepwidestar[1,2,3]{e=e'}{\rChain{7,8}}
\end{tabproofwide}

\FormulaThmDeltaK[Existenz und Eindeutigkeit des neutralen Elements]{%
  \begin{aligned}[t]
    &\MonoidAlg{A}{\star}{e}\vdash{}\\[-2pt]
    &\exists! e'\in A\,\Bigl(
      \forall x\in A\,(e'\star x=x)\land
      \forall x\in A\,(x\star e'=x)
      \Bigr)
  \end{aligned}
}{MonoidIdentityUnique}{%
  \DeltaRow{Mengen}{A\dsep e\dsep e'\dsep x}
}
\begin{tabproofwide}
    \proofstepwidestar[1]{\MonoidAlg{A}{\star}{e}}{\rA}
    \proofstepwidestar[1]{e\in A}{%
      \FormulaRefAuto{MonoidIdentityCarrierAxiom}[ax]{1}}
    \proofstepwidestar[3]{x\in A}{\rA}
    \proofstepwidestar[1,3]{e\star x=x}{%
      \FormulaRefAuto{MonoidLeftIdentityAxiom}[ax]{1,3}}
    \proofstepwidestar[1,3]{x\star e=x}{%
      \FormulaRefAuto{MonoidRightIdentityAxiom}[ax]{1,3}}
    \proofstepwidestar[1]{\forall x\in A\,(e\star x=x)}{%
      \rUI{\rRI{3,4}}}
    \proofstepwidestar[1]{\forall x\in A\,(x\star e=x)}{%
      \rUI{\rRI{3,5}}}
    \proofstepwidestar[1]{%
      \begin{aligned}[t]
        &e\in A\land\Bigl(
          \forall x\in A\,(e\star x=x)\land{}\\[-2pt]
        &\qquad \forall x\in A\,(x\star e=x)
          \Bigr)
      \end{aligned}}{%
      \rAI{2,\rAI{6,7}}}
    \proofstepwidestar[9]{%
      \begin{aligned}[t]
        &e'\in A\land\Bigl(
          \forall x\in A\,(e'\star x=x)\land{}\\[-2pt]
        &\qquad \forall x\in A\,(x\star e'=x)
          \Bigr)
      \end{aligned}}{\rA}
    \proofstepwidestar[9]{e'\in A}{\rAEa{9}}
    \proofstepwidestar[9]{\forall x\in A\,(e'\star x=x)}{%
      \rAEa{\rAEb{9}}}
    \proofstepwidestar[9]{\forall x\in A\,(x\star e'=x)}{%
      \rAEb{\rAEb{9}}}
    \proofstepwidestar[1,9]{e=e'}{%
      \FormulaRefAuto{MonoidIdentityCandidateEqualsIdentity}[thm]{%
        1,10,11,12}}
    \proofstepwidestar[1]{%
      \begin{aligned}[t]
        &\exists! e'\in A\,\Bigl(
          \forall x\in A\,(e'\star x=x)\land{}\\[-2pt]
        &\qquad \forall x\in A\,(x\star e'=x)
          \Bigr)
      \end{aligned}}{\UEI{8,9,13}}
\end{tabproofwide}

\section{Kommutative Monoide}

\FormulaDefDeltaK[Begriff des kommutativen Monoids]{%
  \CommutativeMonoidAlg{A}{\star}{e}%
}{CommutativeMonoidDef}{%
  \DeltaRow{Mengen}{A\dsep e\dsep x\dsep y}
  \DeltaRow{\textbf{Axiome}}{}
  \DeltaPrem{Monoid}{\MonoidAlg{A}{\star}{e}}
  \DeltaRow{Kommutativität}{x,y\in A\vdash x\star y=y\star x}
  \DeltaRow{\textbf{Neues Symbol}}{}
  \DeltaPrem{Kommutatives Monoid}{%
    \CommutativeMonoidAlg{A}{\star}{e}}
}

\FormulaAxiomDeltaK[Zugriff auf das Monoid]{%
  \CommutativeMonoidAlg{A}{\star}{e}
  \vdash\MonoidAlg{A}{\star}{e}%
}{CommutativeMonoidBaseMonoidAxiom}{%
  \DeltaRow{Mengen}{A\dsep e}
}

\FormulaAxiomDeltaK[Zugriff auf die Kommutativität eines Monoids]{%
  \CommutativeMonoidAlg{A}{\star}{e}\dsep x,y\in A
  \vdash x\star y=y\star x%
}{CommutativeMonoidCommutativityAxiom}{%
  \DeltaRow{Mengen}{A\dsep e\dsep x\dsep y}
}

\FormulaThmDeltaK[Kommutative Monoide sind kommutative Halbgruppen]{%
  \CommutativeMonoidAlg{A}{\star}{e}
  \vdash\CommutativeSemigroupAlg{A}{\star}%
}{CommutativeMonoidIsCommutativeSemigroup}{%
  \DeltaRow{Mengen}{A\dsep e\dsep x\dsep y}
}
\begin{tabproof}
  \proofstep{1}{\CommutativeMonoidAlg{A}{\star}{e}}{\rA}
  \proofstep{1}{\MonoidAlg{A}{\star}{e}}{%
    \FormulaRefAuto{CommutativeMonoidBaseMonoidAxiom}[ax]{1}}
  \proofstep{1}{\SemigroupAlg{A}{\star}}{%
    \FormulaRefAuto{MonoidBaseSemigroupAxiom}[ax]{2}}
  \proofstep{4}{x\in A}{\rA}
  \proofstep{5}{y\in A}{\rA}
  \proofstep{1,4,5}{x\star y=y\star x}{%
    \FormulaRefAuto{CommutativeMonoidCommutativityAxiom}[ax]{1,4,5}}
  \proofstep{1}{\CommutativeSemigroupAlg{A}{\star}}{%
    \FormulaRefAuto{CommutativeSemigroupDef}[def]{3,6}}
\end{tabproof}

\section{Einheiten, Teilbarkeit, Irreduzibilität und Primelemente}

\FormulaDefDeltaK[Einheit eines Monoids]{%
  \MonoidUnit{A,\star,e}{u}%
}{MonoidUnitDef}{%
  \DeltaRow{Mengen}{A\dsep e\dsep u\dsep v}
  \DeltaRow{Funktionen}{\star}
  \DeltaPrem{Monoid}{\MonoidAlg{A}{\star}{e}}
  \DeltaRow{Träger}{u\in A}
  \DeltaRow{Inverses Element}{%
    \exists v\in A\,(u\star v=e\land v\star u=e)}
  \DeltaRow{\textbf{Neues Symbol}}{}
  \DeltaPrem{Einheit}{\MonoidUnit{A,\star,e}{u}}
}

\FormulaThmDeltaK[Das neutrale Element ist eine Einheit]{%
  \MonoidAlg{A}{\star}{e}
  \vdash\MonoidUnit{A,\star,e}{e}%
}{MonoidIdentityIsUnit}{%
  \DeltaRow{Mengen}{A\dsep e\dsep v}
  \DeltaRow{Funktionen}{\star}
}
\begin{tabproof}
  \proofstep{1}{\MonoidAlg{A}{\star}{e}}{\rA}
  \proofstep{1}{e\in A}{%
    \FormulaRefAuto{MonoidIdentityCarrierAxiom}[ax]{1}}
  \proofstep{1}{e\star e=e}{%
    \FormulaRefAuto{MonoidLeftIdentityAxiom}[ax]{1,2}}
  \proofstep{1}{%
    \exists v\in A\,(e\star v=e\land v\star e=e)}{%
    \rEI{2,\rAI{3,3}}}
  \proofstep{1}{\MonoidUnit{A,\star,e}{e}}{%
    \FormulaRefAuto{MonoidUnitDef}[def]{1,2,4}}
\end{tabproof}

\FormulaDefDeltaK[Teilbarkeit in einem kommutativen Monoid]{%
  \MonoidDivides{A,\star,e}{a}{b}%
}{CommutativeMonoidDividesDef}{%
  \DeltaRow{Mengen}{A\dsep a\dsep b\dsep c\dsep e}
  \DeltaRow{Funktionen}{\star}
  \DeltaPrem{Kommutatives Monoid}{%
    \CommutativeMonoidAlg{A}{\star}{e}}
  \DeltaRow{Elemente}{a,b\in A}
  \DeltaRow{Quotient}{\exists c\in A\,(b=a\star c)}
  \DeltaRow{\textbf{Neues Symbol}}{}
  \DeltaPrem{Teilbarkeit}{\MonoidDivides{A,\star,e}{a}{b}}
}

\FormulaThmDeltaK[Reflexivität der Teilbarkeit]{%
  \CommutativeMonoidAlg{A}{\star}{e}\dsep a\in A
  \vdash\MonoidDivides{A,\star,e}{a}{a}%
}{CommutativeMonoidDividesReflexive}{%
  \DeltaRow{Mengen}{A\dsep a\dsep c\dsep e}
  \DeltaRow{Funktionen}{\star}
}
\begin{tabproof}
  \proofstep{1}{\CommutativeMonoidAlg{A}{\star}{e}}{\rA}
  \proofstep{2}{a\in A}{\rA}
  \proofstep{1}{\MonoidAlg{A}{\star}{e}}{%
    \FormulaRefAuto{CommutativeMonoidBaseMonoidAxiom}[ax]{1}}
  \proofstep{1}{e\in A}{%
    \FormulaRefAuto{MonoidIdentityCarrierAxiom}[ax]{3}}
  \proofstep{1,2}{a\star e=a}{%
    \FormulaRefAuto{MonoidRightIdentityAxiom}[ax]{3,2}}
  \proofstep{1,2}{a=a\star e}{%
    \FormulaRefAuto{a=b\vdash b=a}{5}}
  \proofstep{1,2}{\exists c\in A\,(a=a\star c)}{%
    \rEI{4,6}}
  \proofstep{1,2}{\MonoidDivides{A,\star,e}{a}{a}}{%
    \FormulaRefAuto{CommutativeMonoidDividesDef}[def]{1,2,2,7}}
\end{tabproof}

\FormulaThmDeltaK[Transitivität der Teilbarkeit]{%
  \begin{aligned}[t]
    &\CommutativeMonoidAlg{A}{\star}{e}\dsep a,b,c\in A\dsep
      \MonoidDivides{A,\star,e}{a}{b}\dsep
      \MonoidDivides{A,\star,e}{b}{c}\vdash{}\\[-2pt]
    &\MonoidDivides{A,\star,e}{a}{c}
  \end{aligned}
}{CommutativeMonoidDividesTransitive}{%
  \DeltaRow{Mengen}{A\dsep a\dsep b\dsep c\dsep e\dsep r\dsep s}
  \DeltaRow{Funktionen}{\star}
}
\begin{tabproofwide}
  \proofstepwidestar[1]{\CommutativeMonoidAlg{A}{\star}{e}}{\rA}
  \proofstepwidestar[2]{a\in A}{\rA}
  \proofstepwidestar[3]{b\in A}{\rA}
  \proofstepwidestar[4]{c\in A}{\rA}
  \proofstepwidestar[5]{\MonoidDivides{A,\star,e}{a}{b}}{\rA}
  \proofstepwidestar[6]{\MonoidDivides{A,\star,e}{b}{c}}{\rA}
  \proofstepwidestar[5]{\exists r\in A\,(b=a\star r)}{%
    \FormulaRefAuto{CommutativeMonoidDividesDef}[def]{5}}
  \proofstepwidestar[6]{\exists s\in A\,(c=b\star s)}{%
    \FormulaRefAuto{CommutativeMonoidDividesDef}[def]{6}}
  \proofstepwidestar[9]{r\in A}{\rA}
  \proofstepwidestar[10]{b=a\star r}{\rA}
  \proofstepwidestar[11]{s\in A}{\rA}
  \proofstepwidestar[12]{c=b\star s}{\rA}
  \proofstepwidestar[1]{\MonoidAlg{A}{\star}{e}}{%
    \FormulaRefAuto{CommutativeMonoidBaseMonoidAxiom}[ax]{1}}
  \proofstepwidestar[1]{\SemigroupAlg{A}{\star}}{%
    \FormulaRefAuto{MonoidBaseSemigroupAxiom}[ax]{13}}
  \proofstepwidestar[1,9,11]{r\star s\in A}{%
    \FormulaRefAuto{SemigroupClosureAxiom}[ax]{14,9,11}}
  \proofstepwide[10,11]{b\star s}{=}{(a\star r)\star s}{\rIE{10}}
  \proofstepwide[1,9,11]{}{=}{a\star(r\star s)}{%
    \FormulaRefAuto{SemigroupAssociativityAxiom}[ax]{14,2,9,11}}
  \proofstepwidestar[1,9,10,11,12]{c=a\star(r\star s)}{%
    \rChain{12,16,17}}
  \proofstepwidestar[1,9,10,11,12]{%
    \exists r\in A\,(c=a\star r)}{\rEI{15,18}}
  \proofstepwidestar[1,2,4,9,10,11,12]{%
    \MonoidDivides{A,\star,e}{a}{c}}{%
    \FormulaRefAuto{CommutativeMonoidDividesDef}[def]{1,2,4,19}}
  \proofstepwidestar[1,2,3,4,5,6]{%
    \MonoidDivides{A,\star,e}{a}{c}}{%
    \rEE{7,9,10,\rEE{8,11,12,20}}}
\end{tabproofwide}

\FormulaThmDeltaK[Einheiten sind genau die Teiler des neutralen Elements]{%
  \begin{aligned}[t]
    &\CommutativeMonoidAlg{A}{\star}{e}\dsep u\in A\vdash{}\\[-2pt]
    &\MonoidUnit{A,\star,e}{u}
      \leftrightarrow\MonoidDivides{A,\star,e}{u}{e}
  \end{aligned}
}{CommutativeMonoidUnitIffDividesIdentity}{%
  \DeltaRow{Mengen}{A\dsep e\dsep u\dsep v\dsep w}
  \DeltaRow{Funktionen}{\star}
}
\begin{tabproofwide}
  \proofstepwidestar[1]{\CommutativeMonoidAlg{A}{\star}{e}}{\rA}
  \proofstepwidestar[2]{u\in A}{\rA}
  \proofstepwidestar[1]{\MonoidAlg{A}{\star}{e}}{%
    \FormulaRefAuto{CommutativeMonoidBaseMonoidAxiom}[ax]{1}}
  \proofstepwidestar[1]{e\in A}{%
    \FormulaRefAuto{MonoidIdentityCarrierAxiom}[ax]{3}}
  \proofstepwidestar[5]{\MonoidUnit{A,\star,e}{u}}{\rA}
  \proofstepwidestar[5]{%
    \exists v\in A\,(u\star v=e\land v\star u=e)}{%
    \FormulaRefAuto{MonoidUnitDef}[def]{5}}
  \proofstepwidestar[7]{v\in A}{\rA}
  \proofstepwidestar[8]{u\star v=e\land v\star u=e}{\rA}
  \proofstepwidestar[8]{u\star v=e}{\rAEa{8}}
  \proofstepwidestar[8]{e=u\star v}{%
    \FormulaRefAuto{a=b\vdash b=a}{9}}
  \proofstepwidestar[7,8]{\exists v\in A\,(e=u\star v)}{%
    \rEI{7,10}}
  \proofstepwidestar[1,2,7,8]{%
    \MonoidDivides{A,\star,e}{u}{e}}{%
    \FormulaRefAuto{CommutativeMonoidDividesDef}[def]{1,2,4,11}}
  \proofstepwidestar[1,2,5]{\MonoidDivides{A,\star,e}{u}{e}}{%
    \rEE{6,7,8,12}}
  \proofstepwidestar[1,2]{%
    \MonoidUnit{A,\star,e}{u}\rightarrow
    \MonoidDivides{A,\star,e}{u}{e}}{\rRI{5,13}}
  \proofstepwidestar[15]{\MonoidDivides{A,\star,e}{u}{e}}{\rA}
  \proofstepwidestar[15]{\exists w\in A\,(e=u\star w)}{%
    \FormulaRefAuto{CommutativeMonoidDividesDef}[def]{15}}
  \proofstepwidestar[17]{w\in A}{\rA}
  \proofstepwidestar[18]{e=u\star w}{\rA}
  \proofstepwidestar[18]{u\star w=e}{%
    \FormulaRefAuto{a=b\vdash b=a}{18}}
  \proofstepwidestar[1,2,17]{w\star u=u\star w}{%
    \FormulaRefAuto{CommutativeMonoidCommutativityAxiom}[ax]{1,17,2}}
  \proofstepwidestar[1,2,17,18]{w\star u=e}{\rChain{20,19}}
  \proofstepwidestar[1,2,17,18]{%
    \exists w\in A\,(u\star w=e\land w\star u=e)}{%
    \rEI{17,\rAI{19,21}}}
  \proofstepwidestar[1,2,17,18]{\MonoidUnit{A,\star,e}{u}}{%
    \FormulaRefAuto{MonoidUnitDef}[def]{3,2,22}}
  \proofstepwidestar[1,2,15]{\MonoidUnit{A,\star,e}{u}}{%
    \rEE{16,17,18,23}}
  \proofstepwidestar[1,2]{%
    \MonoidDivides{A,\star,e}{u}{e}\rightarrow
    \MonoidUnit{A,\star,e}{u}}{\rRI{15,24}}
  \proofstepwidestar[1,2]{%
    \MonoidUnit{A,\star,e}{u}
    \leftrightarrow\MonoidDivides{A,\star,e}{u}{e}}{%
    \rLRI{14,25}}
\end{tabproofwide}

\FormulaDefDeltaK[Irreduzibles Element eines kommutativen Monoids]{%
  \MonoidIrreducible{A,\star,e}{p}%
}{CommutativeMonoidIrreducibleDef}{%
  \DeltaRow{Mengen}{A\dsep a\dsep b\dsep e\dsep p}
  \DeltaRow{Funktionen}{\star}
  \DeltaPrem{Kommutatives Monoid}{%
    \CommutativeMonoidAlg{A}{\star}{e}}
  \DeltaRow{Element}{p\in A}
  \DeltaRow{Keine Einheit}{\neg\MonoidUnit{A,\star,e}{p}}
  \DeltaRow{Irreduzibilität}{%
    \begin{aligned}[t]
      &a,b\in A\dsep p=a\star b\vdash{}\\[-2pt]
      &\MonoidUnit{A,\star,e}{a}\lor
       \MonoidUnit{A,\star,e}{b}
    \end{aligned}}
  \DeltaRow{\textbf{Neues Symbol}}{}
  \DeltaPrem{Irreduzibles Element}{%
    \MonoidIrreducible{A,\star,e}{p}}
}

\FormulaDefDeltaK[Primelement eines kommutativen Monoids]{%
  \MonoidPrime{A,\star,e}{p}%
}{CommutativeMonoidPrimeDef}{%
  \DeltaRow{Mengen}{A\dsep a\dsep b\dsep e\dsep p}
  \DeltaRow{Funktionen}{\star}
  \DeltaPrem{Kommutatives Monoid}{%
    \CommutativeMonoidAlg{A}{\star}{e}}
  \DeltaRow{Element}{p\in A}
  \DeltaRow{Keine Einheit}{\neg\MonoidUnit{A,\star,e}{p}}
  \DeltaRow{Primeigenschaft}{%
    \begin{aligned}[t]
      &a,b\in A\dsep
        \MonoidDivides{A,\star,e}{p}{a\star b}\vdash{}\\[-2pt]
      &\MonoidDivides{A,\star,e}{p}{a}\lor
       \MonoidDivides{A,\star,e}{p}{b}
    \end{aligned}}
  \DeltaRow{\textbf{Neues Symbol}}{}
  \DeltaPrem{Primelement}{\MonoidPrime{A,\star,e}{p}}
}

\FormulaThmDeltaK[Ein geteilter Faktor erzwingt eine Einheit]{%
  \begin{aligned}[t]
    &\CommutativeMonoidAlg{A}{\star}{e}\dsep
      \mathsf{K\ddot{u}rz}\!\left(A,\star\right)\dsep
      p,a,b\in A\dsep p=a\star b\dsep{}\\[-2pt]
    &\MonoidDivides{A,\star,e}{p}{a}
      \vdash\MonoidUnit{A,\star,e}{b}
  \end{aligned}
}{CommutativeCancellativeMonoidFactorDivisorUnit}{%
  \DeltaRow{Mengen}{A\dsep a\dsep b\dsep c\dsep e\dsep p}
  \DeltaRow{Funktionen}{\star}
}
\begin{tabproofwide}
  \proofstepwidestar[1]{\CommutativeMonoidAlg{A}{\star}{e}}{\rA}
  \proofstepwidestar[2]{\mathsf{K\ddot{u}rz}\!\left(A,\star\right)}{\rA}
  \proofstepwidestar[3]{p\in A}{\rA}
  \proofstepwidestar[4]{a\in A}{\rA}
  \proofstepwidestar[5]{b\in A}{\rA}
  \proofstepwidestar[6]{p=a\star b}{\rA}
  \proofstepwidestar[7]{\MonoidDivides{A,\star,e}{p}{a}}{\rA}
  \proofstepwidestar[1]{\MonoidAlg{A}{\star}{e}}{%
    \FormulaRefAuto{CommutativeMonoidBaseMonoidAxiom}[ax]{1}}
  \proofstepwidestar[1]{\SemigroupAlg{A}{\star}}{%
    \FormulaRefAuto{MonoidBaseSemigroupAxiom}[ax]{8}}
  \proofstepwidestar[1]{e\in A}{%
    \FormulaRefAuto{MonoidIdentityCarrierAxiom}[ax]{8}}
  \proofstepwidestar[2]{\mathsf{LK\ddot{u}rz}\!\left(A,\star\right)}{%
    \FormulaRefAuto{CancellativeLeftAxiom}[ax]{2}}
  \proofstepwidestar[7]{\exists c\in A\,(a=p\star c)}{%
    \FormulaRefAuto{CommutativeMonoidDividesDef}[def]{7}}
  \proofstepwidestar[13]{c\in A}{\rA}
  \proofstepwidestar[14]{a=p\star c}{\rA}
  \proofstepwidestar[1,5,13]{c\star b\in A}{%
    \FormulaRefAuto{SemigroupClosureAxiom}[ax]{9,13,5}}
  \proofstepwide[1,3]{p\star e}{=}{p}{%
    \FormulaRefAuto{MonoidRightIdentityAxiom}[ax]{8,3}}
  \proofstepwide[6]{}{=}{a\star b}{\rA}
  \proofstepwide[5,14]{}{=}{(p\star c)\star b}{\rIE{14}}
  \proofstepwide[1,3,5,13]{}{=}{p\star(c\star b)}{%
    \FormulaRefAuto{SemigroupAssociativityAxiom}[ax]{9,3,13,5}}
  \proofstepwidestar[1,3,5,6,13,14]{%
    p\star e=p\star(c\star b)}{\rChain{16,19}}
  \proofstepwidestar[1,2,3,5,6,13,14]{e=c\star b}{%
    \FormulaRefAuto{LeftCancellationAxiom}[ax]{11,3,10,15,20}}
  \proofstepwidestar[1,2,3,5,6,13,14]{c\star b=e}{%
    \FormulaRefAuto{a=b\vdash b=a}{21}}
  \proofstepwidestar[1,5,13]{b\star c=c\star b}{%
    \FormulaRefAuto{CommutativeMonoidCommutativityAxiom}[ax]{1,5,13}}
  \proofstepwidestar[1,2,3,5,6,13,14]{b\star c=e}{%
    \rChain{23,22}}
  \proofstepwidestar[1,2,3,5,6,13,14]{%
    \exists c\in A\,(b\star c=e\land c\star b=e)}{%
    \rEI{13,\rAI{24,22}}}
  \proofstepwidestar[1,2,3,5,6,13,14]{%
    \MonoidUnit{A,\star,e}{b}}{%
    \FormulaRefAuto{MonoidUnitDef}[def]{8,5,25}}
  \proofstepwidestar[1,2,3,4,5,6,7]{%
    \MonoidUnit{A,\star,e}{b}}{\rEE{12,13,14,26}}
\end{tabproofwide}

\FormulaThmDeltaK[Primelemente in kommutativen kürzbaren Monoiden sind irreduzibel]{%
  \begin{aligned}[t]
    &\CommutativeMonoidAlg{A}{\star}{e}\dsep
      \mathsf{K\ddot{u}rz}\!\left(A,\star\right)\dsep
      \MonoidPrime{A,\star,e}{p}\vdash{}\\[-2pt]
    &\MonoidIrreducible{A,\star,e}{p}
  \end{aligned}
}{CommutativeCancellativeMonoidPrimeIsIrreducible}{%
  \DeltaRow{Mengen}{A\dsep a\dsep b\dsep e\dsep p}
  \DeltaRow{Funktionen}{\star}
}
\begin{tabproofwide}
  \proofstepwidestar[1]{\CommutativeMonoidAlg{A}{\star}{e}}{\rA}
  \proofstepwidestar[2]{\mathsf{K\ddot{u}rz}\!\left(A,\star\right)}{\rA}
  \proofstepwidestar[3]{\MonoidPrime{A,\star,e}{p}}{\rA}
  \proofstepwidestar[3]{p\in A}{%
    \FormulaRefAuto{CommutativeMonoidPrimeDef}[def]{3}}
  \proofstepwidestar[3]{\neg\MonoidUnit{A,\star,e}{p}}{%
    \FormulaRefAuto{CommutativeMonoidPrimeDef}[def]{3}}
  \proofstepwidestar[6]{a\in A}{\rA}
  \proofstepwidestar[7]{b\in A}{\rA}
  \proofstepwidestar[8]{p=a\star b}{\rA}
  \proofstepwidestar[1,3]{\MonoidDivides{A,\star,e}{p}{p}}{%
    \FormulaRefAuto{CommutativeMonoidDividesReflexive}[thm]{1,4}}
  \proofstepwidestar[1,3,8]{%
    \MonoidDivides{A,\star,e}{p}{a\star b}}{\rIE{8}}
  \proofstepwidestar[1,3,6,7,8]{%
    \MonoidDivides{A,\star,e}{p}{a}\lor
    \MonoidDivides{A,\star,e}{p}{b}}{%
    \FormulaRefAuto{CommutativeMonoidPrimeDef}[def]{3,6,7,10}}
  \proofstepwidestar[12]{\MonoidDivides{A,\star,e}{p}{a}}{\rA}
  \proofstepwidestar[1,2,3,6,7,8,12]{\MonoidUnit{A,\star,e}{b}}{%
    \FormulaRefAuto{CommutativeCancellativeMonoidFactorDivisorUnit}[thm]{%
      1,2,4,6,7,8,12}}
  \proofstepwidestar[1,2,3,6,7,8,12]{%
    \MonoidUnit{A,\star,e}{a}\lor\MonoidUnit{A,\star,e}{b}}{%
    \rOIb{13}}
  \proofstepwidestar[15]{\MonoidDivides{A,\star,e}{p}{b}}{\rA}
  \proofstepwidestar[1,6,7]{a\star b=b\star a}{%
    \FormulaRefAuto{CommutativeMonoidCommutativityAxiom}[ax]{1,6,7}}
  \proofstepwidestar[1,6,7,8]{p=b\star a}{\rChain{8,16}}
  \proofstepwidestar[1,2,3,6,7,8,15]{\MonoidUnit{A,\star,e}{a}}{%
    \FormulaRefAuto{CommutativeCancellativeMonoidFactorDivisorUnit}[thm]{%
      1,2,4,7,6,17,15}}
  \proofstepwidestar[1,2,3,6,7,8,15]{%
    \MonoidUnit{A,\star,e}{a}\lor\MonoidUnit{A,\star,e}{b}}{%
    \rOIa{18}}
  \proofstepwidestar[1,2,3,6,7,8]{%
    \MonoidUnit{A,\star,e}{a}\lor\MonoidUnit{A,\star,e}{b}}{%
    \rOE{11,12,14,15,19}}
  \proofstepwidestar[1,2,3]{\MonoidIrreducible{A,\star,e}{p}}{%
    \FormulaRefAuto{CommutativeMonoidIrreducibleDef}[def]{%
      1,4,5,6,7,8,20}}
\end{tabproofwide}

\section{Morphismen von Monoiden}

\FormulaDefDeltaK[Begriff des Monoidhomomorphismus]{%
  \MonoidHom{f}{A,\star,e}{B,\diamond,u}%
}{MonoidHomDef}{%
  \DeltaRow{Mengen}{A\dsep B\dsep e\dsep u}
  \DeltaRow{Funktionen}{f\dsep\star\dsep\diamond}
  \DeltaRow{\textbf{Axiome}}{}
  \DeltaPrem{Quellmonoid}{\MonoidAlg{A}{\star}{e}}
  \DeltaPrem{Zielmonoid}{\MonoidAlg{B}{\diamond}{u}}
  \DeltaPrem{Halbgruppenhomomorphismus}{%
    \SemigroupHom{f}{A,\star}{B,\diamond}}
  \DeltaRow{Erhaltung des neutralen Elements}{f(e)=u}
  \DeltaRow{\textbf{Neues Symbol}}{}
  \DeltaPrem{Monoidhomomorphismus}{%
    \MonoidHom{f}{A,\star,e}{B,\diamond,u}}
}

\FormulaAxiomDeltaK[Operationserhaltung eines Monoidhomomorphismus]{%
  \MonoidHom{f}{A,\star,e}{B,\diamond,u}\dsep x,y\in A
  \vdash f(x\star y)=f(x)\diamond f(y)%
}{MonoidHomOperationAxiom}{%
  \DeltaRow{Mengen}{A\dsep B\dsep e\dsep u\dsep x\dsep y}
  \DeltaRow{Funktionen}{f\dsep\star\dsep\diamond}
}

\FormulaAxiomDeltaK[Erhaltung des neutralen Elements]{%
  \MonoidHom{f}{A,\star,e}{B,\diamond,u}
  \vdash f(e)=u%
}{MonoidHomIdentityAxiom}{%
  \DeltaRow{Mengen}{A\dsep B\dsep e\dsep u}
  \DeltaRow{Funktionen}{f\dsep\star\dsep\diamond}
}

\FormulaDefDeltaK[Begriff des Monoidisomorphismus]{%
  \MonoidIso{f}{A,\star,e}{B,\diamond,u}%
}{MonoidIsoDef}{%
  \DeltaRow{Mengen}{A\dsep B\dsep e\dsep u}
  \DeltaRow{Funktionen}{f\dsep\star\dsep\diamond}
  \DeltaRow{\textbf{Axiome}}{}
  \DeltaPrem{Monoidhomomorphismus}{%
    \MonoidHom{f}{A,\star,e}{B,\diamond,u}}
  \DeltaPrem{Bijektivität}{f\colon A\bij B}
  \DeltaRow{\textbf{Neues Symbol}}{}
  \DeltaPrem{Monoidisomorphismus}{%
    \MonoidIso{f}{A,\star,e}{B,\diamond,u}}
}

\subsection{Grundlegende Eigenschaften}

\FormulaThmDeltaK[Komposition von Monoidhomomorphismen]{%
  \begin{aligned}[t]
    &\MonoidHom{f}{A,\star,e}{B,\diamond,u}\dsep
      \MonoidHom{g}{B,\diamond,u}{C,\bullet,v}\\
    &\vdash\MonoidHom{g\circ f}{A,\star,e}{C,\bullet,v}
  \end{aligned}
}{MonoidHomComposition}{%
  \DeltaRow{Mengen}{A\dsep B\dsep C\dsep e\dsep u\dsep v}
  \DeltaRow{Funktionen}{f\dsep g\dsep\star\dsep\diamond\dsep\bullet}
}
\begin{tabproofwide}
  \proofstepwidestar[1]{\MonoidHom{f}{A,\star,e}{B,\diamond,u}}{\rA}
  \proofstepwidestar[2]{\MonoidHom{g}{B,\diamond,u}{C,\bullet,v}}{\rA}
  \proofstepwidestar[1]{\SemigroupHom{f}{A,\star}{B,\diamond}}{%
    \FormulaRefAuto{MonoidHomDef}[def]{1}}
  \proofstepwidestar[2]{\SemigroupHom{g}{B,\diamond}{C,\bullet}}{%
    \FormulaRefAuto{MonoidHomDef}[def]{2}}
  \proofstepwidestar[1,2]{%
    \SemigroupHom{g\circ f}{A,\star}{C,\bullet}}{%
    \FormulaRefAuto{SemigroupHomComposition}[thm]{3,4}}
  \proofstepwidestar[1]{f(e)=u}{%
    \FormulaRefAuto{MonoidHomIdentityAxiom}[ax]{1}}
  \proofstepwidestar[2]{g(u)=v}{%
    \FormulaRefAuto{MonoidHomIdentityAxiom}[ax]{2}}
  \proofstepwidestar[1]{\MonoidAlg{A}{\star}{e}}{%
    \FormulaRefAuto{MonoidHomDef}[def]{1}}
  \proofstepwidestar[1]{e\in A}{%
    \FormulaRefAuto{MonoidIdentityCarrierAxiom}[ax]{8}}
  \proofstepwidestar[1]{f\colon A\to B}{%
    \FormulaRefAuto{SemigroupHomFunctionAxiom}[ax]{3}}
  \proofstepwidestar[2]{g\colon B\to C}{%
    \FormulaRefAuto{SemigroupHomFunctionAxiom}[ax]{4}}
  \proofstepwide[1,2]{(g\circ f)(e)}{=}{g(f(e))}{%
    \FormulaRefAuto{x\in A \vdash (G\circ F)(x)=G(F(x))}{%
      10,11,9}}
  \proofstepwide[1,2]{}{=}{g(u)}{\rIE{6}}
  \proofstepwide[1,2]{}{=}{v}{\rIE{7}}
  \proofstepwidestar[1,2]{(g\circ f)(e)=v}{\rChain{12,14}}
  \proofstepwidestar[2]{\MonoidAlg{C}{\bullet}{v}}{%
    \FormulaRefAuto{MonoidHomDef}[def]{2}}
  \proofstepwidestar[1,2]{%
    \MonoidHom{g\circ f}{A,\star,e}{C,\bullet,v}}{%
    \FormulaRefAuto{MonoidHomDef}[def]{8,16,5,15}}
\end{tabproofwide}

\subsection{Kommutative Homomorphismen und Isomorphismen}

\FormulaDefDeltaK[Homomorphismen kommutativer Monoide]{%
  \begin{aligned}[t]
    &\CommutativeMonoidHom{f}{A,\star,e}{B,\diamond,u}\\[-2pt]
    &\quad\coloneqq
      \CommutativeMonoidAlg{A}{\star}{e}\land
      \CommutativeMonoidAlg{B}{\diamond}{u}\\[-2pt]
    &\qquad\land\MonoidHom{f}{A,\star,e}{B,\diamond,u}
  \end{aligned}
}{CommutativeMonoidHomDef}{%
  \DeltaRow{Mengen}{A\dsep B\dsep e\dsep u}
  \DeltaRow{Funktionen}{f\dsep\star\dsep\diamond}
  \DeltaRow{Operationserhaltung}{f(x\star y)=f(x)\diamond f(y)}
  \DeltaRow{Erhaltung des Neutralen}{f(e)=u}
  \DeltaRow{\textbf{Neues Symbol}}{}
  \DeltaPrem{Symbol}{\mathsf{KMonHom}}
}

\FormulaAxiomDeltaK[Quellstruktur eines Homomorphismus kommutativer Monoide]{%
  \CommutativeMonoidHom{f}{A,\star,e}{B,\diamond,u}
  \vdash\CommutativeMonoidAlg{A}{\star}{e}%
}{CommutativeMonoidHomSourceAxiom}{%
  \DeltaRow{Mengen}{A\dsep B\dsep e\dsep u}
  \DeltaRow{Funktionen}{f\dsep\star\dsep\diamond}
}

\FormulaAxiomDeltaK[Zielstruktur eines Homomorphismus kommutativer Monoide]{%
  \CommutativeMonoidHom{f}{A,\star,e}{B,\diamond,u}
  \vdash\CommutativeMonoidAlg{B}{\diamond}{u}%
}{CommutativeMonoidHomTargetAxiom}{%
  \DeltaRow{Mengen}{A\dsep B\dsep e\dsep u}
  \DeltaRow{Funktionen}{f\dsep\star\dsep\diamond}
}

\FormulaAxiomDeltaK[Monoidhomomorphismus eines kommutativen Homomorphismus]{%
  \CommutativeMonoidHom{f}{A,\star,e}{B,\diamond,u}
  \vdash\MonoidHom{f}{A,\star,e}{B,\diamond,u}%
}{CommutativeMonoidHomBaseHomAxiom}{%
  \DeltaRow{Mengen}{A\dsep B\dsep e\dsep u}
  \DeltaRow{Funktionen}{f\dsep\star\dsep\diamond}
}

\FormulaDefDeltaK[Isomorphismen kommutativer Monoide]{%
  \begin{aligned}[t]
    &\CommutativeMonoidIso{f}{A,\star,e}{B,\diamond,u}\\[-2pt]
    &\quad\coloneqq
      \CommutativeMonoidHom{f}{A,\star,e}{B,\diamond,u}
      \land f\colon A\bij B
  \end{aligned}
}{CommutativeMonoidIsoDef}{%
  \DeltaRow{Mengen}{A\dsep B\dsep e\dsep u}
  \DeltaRow{Funktionen}{f\dsep\star\dsep\diamond}
  \DeltaRow{\textbf{Neues Symbol}}{}
  \DeltaPrem{Isomorphismus kommutativer Monoide}{\mathsf{KMonIso}}
}

\FormulaAxiomDeltaK[Homomorphismus eines Isomorphismus kommutativer Monoide]{%
  \CommutativeMonoidIso{f}{A,\star,e}{B,\diamond,u}
  \vdash\CommutativeMonoidHom{f}{A,\star,e}{B,\diamond,u}%
}{CommutativeMonoidIsoHomAxiom}{%
  \DeltaRow{Mengen}{A\dsep B\dsep e\dsep u}
  \DeltaRow{Funktionen}{f\dsep\star\dsep\diamond}
}

\FormulaAxiomDeltaK[Bijektivität eines Isomorphismus kommutativer Monoide]{%
  \CommutativeMonoidIso{f}{A,\star,e}{B,\diamond,u}
  \vdash f\colon A\bij B%
}{CommutativeMonoidIsoBijectiveAxiom}{%
  \DeltaRow{Mengen}{A\dsep B\dsep e\dsep u}
  \DeltaRow{Funktionen}{f\dsep\star\dsep\diamond}
}

\chapter{Halbringe mit Eins}

\section{Definition des Halbrings mit Eins}

\FormulaDefDeltaK[Begriff des Halbrings mit Eins]{%
  \SemiringAlg{R}{+}{\cdot}{0}{1}%
}{SemiringDef}{%
  \DeltaRow{Mengen}{R\dsep 0\dsep 1\dsep a\dsep b\dsep c}
  \DeltaRow{\textbf{Axiome}}{}
  \DeltaPrem{\makecell[l]{Additives kommutatives\\ Monoid}}{%
    \CommutativeMonoidAlg{R}{+}{0}}
  \DeltaPrem{Multiplikatives Monoid}{%
    \MonoidAlg{R}{\cdot}{1}}
  \DeltaRow{Linksdistributivität}{%
    \begin{aligned}[t]
      &a,b,c\in R\vdash{}\\[-2pt]
      &a\cdot(b+c)=a\cdot b+a\cdot c
    \end{aligned}}
  \DeltaRow{Rechtsdistributivität}{%
    \begin{aligned}[t]
      &a,b,c\in R\vdash{}\\[-2pt]
      &(a+b)\cdot c=a\cdot c+b\cdot c
    \end{aligned}}
  \DeltaRow{Linksabsorption}{a\in R\vdash0\cdot a=0}
  \DeltaRow{Rechtsabsorption}{a\in R\vdash a\cdot0=0}
  \DeltaRow{\textbf{Neues Symbol}}{}
  \DeltaPrem{Halbring mit Eins}{%
    \SemiringAlg{R}{+}{\cdot}{0}{1}}
}

\FormulaAxiomDeltaK[Zugriff auf das additive kommutative Monoid]{%
  \SemiringAlg{R}{+}{\cdot}{0}{1}
  \vdash\CommutativeMonoidAlg{R}{+}{0}%
}{SemiringAdditiveCommutativeMonoidAxiom}{%
  \DeltaRow{Mengen}{R\dsep 0\dsep 1}
}

\FormulaAxiomDeltaK[Zugriff auf das multiplikative Monoid]{%
  \SemiringAlg{R}{+}{\cdot}{0}{1}
  \vdash\MonoidAlg{R}{\cdot}{1}%
}{SemiringMultiplicativeMonoidAxiom}{%
  \DeltaRow{Mengen}{R\dsep 0\dsep 1}
}

\FormulaAxiomDeltaK[Zugriff auf die Linksdistributivität]{%
  \SemiringAlg{R}{+}{\cdot}{0}{1}\dsep a,b,c\in R
  \vdash a\cdot(b+c)=a\cdot b+a\cdot c%
}{SemiringLeftDistributiveAxiom}{%
  \DeltaRow{Mengen}{R\dsep 0\dsep 1\dsep a\dsep b\dsep c}
}

\FormulaAxiomDeltaK[Zugriff auf die Rechtsdistributivität]{%
  \SemiringAlg{R}{+}{\cdot}{0}{1}\dsep a,b,c\in R
  \vdash (a+b)\cdot c=a\cdot c+b\cdot c%
}{SemiringRightDistributiveAxiom}{%
  \DeltaRow{Mengen}{R\dsep 0\dsep 1\dsep a\dsep b\dsep c}
}

\FormulaAxiomDeltaK[Zugriff auf die Linksabsorption]{%
  \SemiringAlg{R}{+}{\cdot}{0}{1}\dsep a\in R
  \vdash0\cdot a=0%
}{SemiringLeftZeroAxiom}{%
  \DeltaRow{Mengen}{R\dsep 0\dsep 1\dsep a}
}

\FormulaAxiomDeltaK[Zugriff auf die Rechtsabsorption]{%
  \SemiringAlg{R}{+}{\cdot}{0}{1}\dsep a\in R
  \vdash a\cdot0=0%
}{SemiringRightZeroAxiom}{%
  \DeltaRow{Mengen}{R\dsep 0\dsep 1\dsep a}
}

\section{Kommutative Halbringe mit Eins}

\FormulaDefDeltaK[Kommutativer Halbring mit Eins]{%
  \CommutativeSemiringAlg{R}{+}{\cdot}{0}{1}%
}{CommutativeUnitalSemiringDef}{%
  \DeltaRow{Mengen}{R\dsep 0\dsep 1}
  \DeltaRow{\textbf{Axiome}}{}
  \DeltaPrem{Halbring mit Eins}{%
    \SemiringAlg{R}{+}{\cdot}{0}{1}}
  \DeltaPrem{Kommutatives multiplikatives Monoid}{%
    \CommutativeMonoidAlg{R}{\cdot}{1}}
  \DeltaRow{\textbf{Neues Symbol}}{}
  \DeltaPrem{Kommutativer Halbring mit Eins}{%
    \CommutativeSemiringAlg{R}{+}{\cdot}{0}{1}}
}

\FormulaAxiomDeltaK[Zugriff auf den Halbring mit Eins]{%
  \CommutativeSemiringAlg{R}{+}{\cdot}{0}{1}
  \vdash\SemiringAlg{R}{+}{\cdot}{0}{1}%
}{CommutativeSemiringBaseAxiom}{%
  \DeltaRow{Mengen}{R\dsep 0\dsep 1}
}

\FormulaAxiomDeltaK[Zugriff auf das kommutative multiplikative Monoid]{%
  \CommutativeSemiringAlg{R}{+}{\cdot}{0}{1}
  \vdash\CommutativeMonoidAlg{R}{\cdot}{1}%
}{CommutativeSemiringMultiplicativeMonoidAxiom}{%
  \DeltaRow{Mengen}{R\dsep 0\dsep 1}
}

\section{Morphismen von Halbringen}

\FormulaDefDeltaK[Begriff des Halbringhomomorphismus]{%
  \SemiringHom{f}{R,\oplus,\odot,0_R,1_R}
    {S,\boxplus,\boxtimes,0_S,1_S}%
}{SemiringHomDef}{%
  \DeltaRow{Mengen}{R\dsep S\dsep0_R\dsep1_R\dsep0_S\dsep1_S}
  \DeltaRow{Funktionen}{f\dsep\oplus\dsep\odot\dsep\boxplus\dsep\boxtimes}
  \DeltaRow{\textbf{Axiome}}{}
  \DeltaPrem{Quellhalbring}{%
    \SemiringAlg{R}{\oplus}{\odot}{0_R}{1_R}}
  \DeltaPrem{Zielhalbring}{%
    \SemiringAlg{S}{\boxplus}{\boxtimes}{0_S}{1_S}}
  \DeltaPrem{Additiver Anteil}{%
    \MonoidHom{f}{R,\oplus,0_R}{S,\boxplus,0_S}}
  \DeltaPrem{Multiplikativer Anteil}{%
    \MonoidHom{f}{R,\odot,1_R}{S,\boxtimes,1_S}}
  \DeltaRow{\textbf{Neues Symbol}}{}
  \DeltaPrem{Symbol}{\mathsf{HringHom}}
}

\FormulaAxiomDeltaK[Additionserhaltung eines Halbringhomomorphismus]{%
  \begin{aligned}[t]
    &\SemiringHom{f}{R,\oplus,\odot,0_R,1_R}
      {S,\boxplus,\boxtimes,0_S,1_S}\dsep x,y\in R\\
    &\vdash f(x\oplus y)=f(x)\boxplus f(y)
  \end{aligned}
}{SemiringHomAdditionAxiom}{%
  \DeltaRow{Mengen}{R\dsep S\dsep x\dsep y}
  \DeltaRow{Funktionen}{f\dsep\oplus\dsep\odot\dsep\boxplus\dsep\boxtimes}
}

\FormulaAxiomDeltaK[Trägerabbildung eines Halbringhomomorphismus]{%
  \SemiringHom{f}{R,\oplus,\odot,0_R,1_R}
    {S,\boxplus,\boxtimes,0_S,1_S}
  \vdash f\colon R\to S%
}{SemiringHomFunctionAxiom}{%
  \DeltaRow{Mengen}{R\dsep S}
  \DeltaRow{Funktionen}{f\dsep\oplus\dsep\odot\dsep\boxplus\dsep\boxtimes}
}

\FormulaAxiomDeltaK[Multiplikationserhaltung eines Halbringhomomorphismus]{%
  \begin{aligned}[t]
    &\SemiringHom{f}{R,\oplus,\odot,0_R,1_R}
      {S,\boxplus,\boxtimes,0_S,1_S}\dsep x,y\in R\\
    &\vdash f(x\odot y)=f(x)\boxtimes f(y)
  \end{aligned}
}{SemiringHomMultiplicationAxiom}{%
  \DeltaRow{Mengen}{R\dsep S\dsep x\dsep y}
  \DeltaRow{Funktionen}{f\dsep\oplus\dsep\odot\dsep\boxplus\dsep\boxtimes}
}

\FormulaAxiomDeltaK[Null- und Einserhaltung eines Halbringhomomorphismus]{%
  \SemiringHom{f}{R,\oplus,\odot,0_R,1_R}
    {S,\boxplus,\boxtimes,0_S,1_S}
  \vdash f(0_R)=0_S\land f(1_R)=1_S%
}{SemiringHomConstantsAxiom}{%
  \DeltaRow{Mengen}{R\dsep S\dsep0_R\dsep1_R\dsep0_S\dsep1_S}
  \DeltaRow{Funktionen}{f\dsep\oplus\dsep\odot\dsep\boxplus\dsep\boxtimes}
}

\FormulaDefDeltaK[Begriff des Halbringisomorphismus]{%
  \SemiringIso{f}{R,\oplus,\odot,0_R,1_R}
    {S,\boxplus,\boxtimes,0_S,1_S}%
}{SemiringIsoDef}{%
  \DeltaRow{Mengen}{R\dsep S\dsep0_R\dsep1_R\dsep0_S\dsep1_S}
  \DeltaRow{Funktionen}{f\dsep\oplus\dsep\odot\dsep\boxplus\dsep\boxtimes}
  \DeltaRow{\textbf{Axiome}}{}
  \DeltaPrem{Homomorphismus}{%
    \SemiringHom{f}{R,\oplus,\odot,0_R,1_R}
      {S,\boxplus,\boxtimes,0_S,1_S}}
  \DeltaPrem{Bijektivität}{f\colon R\bij S}
  \DeltaRow{\textbf{Neues Symbol}}{}
  \DeltaPrem{Symbol}{\mathsf{HringIso}}
}

\subsection{Grundlegende Eigenschaften}

\FormulaThmDeltaK[Komposition von Halbringhomomorphismen]{%
  \begin{aligned}[t]
    &\SemiringHom{f}{R,\oplus,\odot,0_R,1_R}
      {S,\boxplus,\boxtimes,0_S,1_S}\dsep{}\\[-2pt]
    &\SemiringHom{g}{S,\boxplus,\boxtimes,0_S,1_S}
      {T,\uplus,\otimes,0_T,1_T}\\
    &\vdash
      \SemiringHom{g\circ f}{R,\oplus,\odot,0_R,1_R}
      {T,\uplus,\otimes,0_T,1_T}
  \end{aligned}
}{SemiringHomComposition}{%
  \DeltaRow{Mengen}{R\dsep S\dsep T}
  \DeltaRow{Funktionen}{f\dsep g\dsep\oplus\dsep\odot\dsep
    \boxplus\dsep\boxtimes\dsep\uplus\dsep\otimes}
}
\begin{tabproofwide}
  \proofstepwidestar[1]{%
    \SemiringHom{f}{R,\oplus,\odot,0_R,1_R}
      {S,\boxplus,\boxtimes,0_S,1_S}}{\rA}
  \proofstepwidestar[2]{%
    \SemiringHom{g}{S,\boxplus,\boxtimes,0_S,1_S}
      {T,\uplus,\otimes,0_T,1_T}}{\rA}
  \proofstepwidestar[1]{%
    \MonoidHom{f}{R,\oplus,0_R}{S,\boxplus,0_S}}{%
    \FormulaRefAuto{SemiringHomDef}[def]{1}}
  \proofstepwidestar[2]{%
    \MonoidHom{g}{S,\boxplus,0_S}{T,\uplus,0_T}}{%
    \FormulaRefAuto{SemiringHomDef}[def]{2}}
  \proofstepwidestar[1,2]{%
    \MonoidHom{g\circ f}{R,\oplus,0_R}{T,\uplus,0_T}}{%
    \FormulaRefAuto{MonoidHomComposition}[thm]{3,4}}
  \proofstepwidestar[1]{%
    \MonoidHom{f}{R,\odot,1_R}{S,\boxtimes,1_S}}{%
    \FormulaRefAuto{SemiringHomDef}[def]{1}}
  \proofstepwidestar[2]{%
    \MonoidHom{g}{S,\boxtimes,1_S}{T,\otimes,1_T}}{%
    \FormulaRefAuto{SemiringHomDef}[def]{2}}
  \proofstepwidestar[1,2]{%
    \MonoidHom{g\circ f}{R,\odot,1_R}{T,\otimes,1_T}}{%
    \FormulaRefAuto{MonoidHomComposition}[thm]{6,7}}
  \proofstepwidestar[1]{%
    \SemiringAlg{R}{\oplus}{\odot}{0_R}{1_R}}{%
    \FormulaRefAuto{SemiringHomDef}[def]{1}}
  \proofstepwidestar[2]{%
    \SemiringAlg{T}{\uplus}{\otimes}{0_T}{1_T}}{%
    \FormulaRefAuto{SemiringHomDef}[def]{2}}
  \proofstepwidestar[1,2]{%
    \SemiringHom{g\circ f}{R,\oplus,\odot,0_R,1_R}
      {T,\uplus,\otimes,0_T,1_T}}{%
    \FormulaRefAuto{SemiringHomDef}[def]{9,10,5,8}}
\end{tabproofwide}

\subsection{Kommutative Homomorphismen und Isomorphismen}

\FormulaDefDeltaK[Homomorphismen kommutativer Halbringe]{%
  \begin{aligned}[t]
    &\CommutativeSemiringHom{f}{R,\oplus,\odot,0_R,1_R}
      {S,\boxplus,\boxtimes,0_S,1_S}\\
    &\quad\coloneqq
      \CommutativeSemiringAlg{R}{\oplus}{\odot}{0_R}{1_R}\land
      \CommutativeSemiringAlg{S}{\boxplus}{\boxtimes}{0_S}{1_S}\\[-2pt]
    &\qquad\land
      \SemiringHom{f}{R,\oplus,\odot,0_R,1_R}
        {S,\boxplus,\boxtimes,0_S,1_S}
  \end{aligned}
}{CommutativeSemiringHomDef}{%
  \DeltaRow{Mengen}{R\dsep S\dsep0_R\dsep1_R\dsep0_S\dsep1_S}
  \DeltaRow{Funktionen}{f\dsep\oplus\dsep\odot\dsep\boxplus\dsep\boxtimes}
  \DeltaRow{\textbf{Neues Symbol}}{}
  \DeltaPrem{\makecell[l]{Homomorphismus\\ kommutativer Halbringe}}{%
    \mathsf{KHringHom}}
}

\FormulaAxiomDeltaK[Quellstruktur eines Homomorphismus kommutativer Halbringe]{%
  \begin{aligned}[t]
    &\CommutativeSemiringHom{f}{R,\oplus,\odot,0_R,1_R}
      {S,\boxplus,\boxtimes,0_S,1_S}\\[-2pt]
    &\vdash\CommutativeSemiringAlg{R}{\oplus}{\odot}{0_R}{1_R}
  \end{aligned}%
}{CommutativeSemiringHomSourceAxiom}{%
  \DeltaRow{Mengen}{R\dsep S\dsep0_R\dsep1_R\dsep0_S\dsep1_S}
  \DeltaRow{Funktionen}{f\dsep\oplus\dsep\odot\dsep\boxplus\dsep\boxtimes}
}

\FormulaAxiomDeltaK[Zielstruktur eines Homomorphismus kommutativer Halbringe]{%
  \begin{aligned}[t]
    &\CommutativeSemiringHom{f}{R,\oplus,\odot,0_R,1_R}
      {S,\boxplus,\boxtimes,0_S,1_S}\\[-2pt]
    &\vdash\CommutativeSemiringAlg{S}{\boxplus}{\boxtimes}{0_S}{1_S}
  \end{aligned}%
}{CommutativeSemiringHomTargetAxiom}{%
  \DeltaRow{Mengen}{R\dsep S\dsep0_R\dsep1_R\dsep0_S\dsep1_S}
  \DeltaRow{Funktionen}{f\dsep\oplus\dsep\odot\dsep\boxplus\dsep\boxtimes}
}

\FormulaAxiomDeltaK[Halbringhomomorphismus eines kommutativen Homomorphismus]{%
  \begin{aligned}[t]
    &\CommutativeSemiringHom{f}{R,\oplus,\odot,0_R,1_R}
      {S,\boxplus,\boxtimes,0_S,1_S}\\[-2pt]
    &\vdash
      \SemiringHom{f}{R,\oplus,\odot,0_R,1_R}
        {S,\boxplus,\boxtimes,0_S,1_S}
  \end{aligned}%
}{CommutativeSemiringHomBaseHomAxiom}{%
  \DeltaRow{Mengen}{R\dsep S\dsep0_R\dsep1_R\dsep0_S\dsep1_S}
  \DeltaRow{Funktionen}{f\dsep\oplus\dsep\odot\dsep\boxplus\dsep\boxtimes}
}

\FormulaDefDeltaK[Isomorphismen kommutativer Halbringe]{%
  \begin{aligned}[t]
    &\CommutativeSemiringIso{f}{R,\oplus,\odot,0_R,1_R}
      {S,\boxplus,\boxtimes,0_S,1_S}\\
    &\quad\coloneqq
      \CommutativeSemiringHom{f}{R,\oplus,\odot,0_R,1_R}
        {S,\boxplus,\boxtimes,0_S,1_S}\\[-2pt]
    &\qquad\land f\colon R\bij S
  \end{aligned}
}{CommutativeSemiringIsoDef}{%
  \DeltaRow{Mengen}{R\dsep S\dsep0_R\dsep1_R\dsep0_S\dsep1_S}
  \DeltaRow{Funktionen}{f\dsep\oplus\dsep\odot\dsep\boxplus\dsep\boxtimes}
  \DeltaRow{\textbf{Neues Symbol}}{}
  \DeltaPrem{Isomorphismus kommutativer Halbringe}{\mathsf{KHringIso}}
}

\FormulaAxiomDeltaK[Homomorphismus eines Isomorphismus kommutativer Halbringe]{%
  \begin{aligned}[t]
    &\CommutativeSemiringIso{f}{R,\oplus,\odot,0_R,1_R}
      {S,\boxplus,\boxtimes,0_S,1_S}\\[-2pt]
    &\vdash
      \CommutativeSemiringHom{f}{R,\oplus,\odot,0_R,1_R}
        {S,\boxplus,\boxtimes,0_S,1_S}
  \end{aligned}%
}{CommutativeSemiringIsoHomAxiom}{%
  \DeltaRow{Mengen}{R\dsep S\dsep0_R\dsep1_R\dsep0_S\dsep1_S}
  \DeltaRow{Funktionen}{f\dsep\oplus\dsep\odot\dsep\boxplus\dsep\boxtimes}
}

\FormulaAxiomDeltaK[Bijektivität eines Isomorphismus kommutativer Halbringe]{%
  \begin{aligned}[t]
    &\CommutativeSemiringIso{f}{R,\oplus,\odot,0_R,1_R}
      {S,\boxplus,\boxtimes,0_S,1_S}\\[-2pt]
    &\vdash f\colon R\bij S
  \end{aligned}%
}{CommutativeSemiringIsoBijectiveAxiom}{%
  \DeltaRow{Mengen}{R\dsep S\dsep0_R\dsep1_R\dsep0_S\dsep1_S}
  \DeltaRow{Funktionen}{f\dsep\oplus\dsep\odot\dsep\boxplus\dsep\boxtimes}
}

\chapter{Die natürlichen Zahlen als Halbring}

\section{Das additive Monoid}

\FormulaThmDeltaK[Die Addition bildet eine Halbgruppe]{%
  \SemigroupAlg{\mathbb{N}}{+}%
}{NaturalAdditionSemigroup}{%
  \DeltaRow{Mengen}{\mathbb{N}\dsep a\dsep b\dsep c}
}
\begin{tabproof}
  \proofstep{1}{a\in\mathbb{N}}{\rA}
  \proofstep{2}{b\in\mathbb{N}}{\rA}
  \proofstep{1,2}{a+b\in\mathbb{N}}{%
    \FormulaRefAuto{PeanoAddClosure}[thm]{1,2}}
  \proofstep{4}{c\in\mathbb{N}}{\rA}
  \proofstep{1,2,4}{(a+b)+c=a+(b+c)}{%
    \FormulaRefAuto{PeanoAddAssociative}[thm]{1,2,4}}
  \proofstep{}{\SemigroupAlg{\mathbb{N}}{+}}{%
    \FormulaRefAuto{SemigroupDef}[def]{3,5}}
\end{tabproof}

\FormulaThmDeltaK[Die Addition bildet ein Monoid]{%
  \MonoidAlg{\mathbb{N}}{+}{0}%
}{NaturalAdditionMonoid}{%
  \DeltaRow{Mengen}{\mathbb{N}\dsep 0\dsep a}
}
\begin{tabproof}
  \proofstep{}{\SemigroupAlg{\mathbb{N}}{+}}{%
    \FormulaRefAuto{NaturalAdditionSemigroup}[thm]}
  \proofstep{}{0\in\mathbb{N}}{%
    \FormulaRefAuto{PeanoZero}[ax]}
  \proofstep{3}{a\in\mathbb{N}}{\rA}
  \proofstep{3}{0+a=a}{%
    \FormulaRefAuto{PeanoAddLeftZero}[thm]{3}}
  \proofstep{3}{a+0=a}{%
    \FormulaRefAuto{PeanoAddRightZero}[thm]{3}}
  \proofstep{}{\MonoidAlg{\mathbb{N}}{+}{0}}{%
    \FormulaRefAuto{MonoidDef}[def]{1,2,4,5}}
\end{tabproof}

\FormulaThmDeltaK[Die Addition bildet ein kommutatives Monoid]{%
  \CommutativeMonoidAlg{\mathbb{N}}{+}{0}%
}{NaturalAdditiveCommutativeMonoid}{%
  \DeltaRow{Mengen}{\mathbb{N}\dsep 0\dsep a\dsep b}
}
\begin{tabproof}
  \proofstep{}{\MonoidAlg{\mathbb{N}}{+}{0}}{%
    \FormulaRefAuto{NaturalAdditionMonoid}[thm]}
  \proofstep{2}{a\in\mathbb{N}}{\rA}
  \proofstep{3}{b\in\mathbb{N}}{\rA}
  \proofstep{2,3}{a+b=b+a}{%
    \FormulaRefAuto{PeanoAddCommutative}[thm]{2,3}}
  \proofstep{}{\CommutativeMonoidAlg{\mathbb{N}}{+}{0}}{%
    \FormulaRefAuto{CommutativeMonoidDef}[def]{1,4}}
\end{tabproof}

\section{Kürzbarkeit der Addition}

\FormulaThmDeltaK[Die Addition ist linkskürzbar]{%
  \mathsf{LK\ddot{u}rz}\!\left(\mathbb{N},+\right)%
}{NaturalAdditionLeftCancellative}{%
  \DeltaRow{Mengen}{\mathbb{N}\dsep a\dsep b\dsep c}
}
\begin{tabproof}
  \proofstep{1}{a\in\mathbb{N}}{\rA}
  \proofstep{2}{b\in\mathbb{N}}{\rA}
  \proofstep{3}{c\in\mathbb{N}}{\rA}
  \proofstep{4}{a+b=a+c}{\rA}
  \proofstep{1,2,3,4}{b=c}{%
    \FormulaRefAuto{PeanoAddLeftCancellation}[thm]{1,2,3,4}}
  \proofstep{}{%
    \mathsf{LK\ddot{u}rz}\!\left(\mathbb{N},+\right)%
  }{\FormulaRefAuto{LeftCancellativeDef}[def]{5}}
\end{tabproof}

\FormulaThmDeltaK[Die Addition ist rechtskürzbar]{%
  \mathsf{RK\ddot{u}rz}\!\left(\mathbb{N},+\right)%
}{NaturalAdditionRightCancellative}{%
  \DeltaRow{Mengen}{\mathbb{N}\dsep a\dsep b\dsep c}
}
\begin{tabproof}
  \proofstep{}{%
    \CommutativeMonoidAlg{\mathbb{N}}{+}{0}}{%
    \FormulaRefAuto{NaturalAdditiveCommutativeMonoid}[thm]}
  \proofstep{}{%
    \CommutativeSemigroupAlg{\mathbb{N}}{+}}{%
    \FormulaRefAuto{CommutativeMonoidIsCommutativeSemigroup}[thm]{1}}
  \proofstep{}{%
    \mathsf{LK\ddot{u}rz}\!\left(\mathbb{N},+\right)}{%
    \FormulaRefAuto{NaturalAdditionLeftCancellative}[thm]}
  \proofstep{}{%
    \mathsf{RK\ddot{u}rz}\!\left(\mathbb{N},+\right)}{%
    \FormulaRefAuto{%
      CommutativeLeftCancellativeImpliesRightCancellative%
    }[thm]{2,3}}
\end{tabproof}

\FormulaThmDeltaK[Die Addition ist kürzbar]{%
  \mathsf{K\ddot{u}rz}\!\left(\mathbb{N},+\right)%
}{NaturalAdditionCancellative}{%
  \DeltaRow{Mengen}{\mathbb{N}\dsep a\dsep b\dsep c}
}
\begin{tabproof}
  \proofstep{}{%
    \mathsf{LK\ddot{u}rz}\!\left(\mathbb{N},+\right)}{%
    \FormulaRefAuto{NaturalAdditionLeftCancellative}[thm]}
  \proofstep{}{%
    \mathsf{RK\ddot{u}rz}\!\left(\mathbb{N},+\right)}{%
    \FormulaRefAuto{NaturalAdditionRightCancellative}[thm]}
  \proofstep{}{%
    \mathsf{K\ddot{u}rz}\!\left(\mathbb{N},+\right)}{%
    \FormulaRefAuto{CancellativeOperationDef}[def]{1,2}}
\end{tabproof}

\section{Das multiplikative Monoid}

\FormulaThmDeltaK[Die Multiplikation bildet eine Halbgruppe]{%
  \SemigroupAlg{\mathbb{N}}{\cdot}%
}{NaturalMultiplicationSemigroup}{%
  \DeltaRow{Mengen}{\mathbb{N}\dsep a\dsep b\dsep c}
}
\begin{tabproof}
  \proofstep{1}{a\in\mathbb{N}}{\rA}
  \proofstep{2}{b\in\mathbb{N}}{\rA}
  \proofstep{1,2}{a\cdot b\in\mathbb{N}}{%
    \FormulaRefAuto{PeanoMulClosure}[thm]{1,2}}
  \proofstep{4}{c\in\mathbb{N}}{\rA}
  \proofstep{1,2,4}{(a\cdot b)\cdot c=a\cdot(b\cdot c)}{%
    \FormulaRefAuto{PeanoMulAssociative}[thm]{1,2,4}}
  \proofstep{}{\SemigroupAlg{\mathbb{N}}{\cdot}}{%
    \FormulaRefAuto{SemigroupDef}[def]{3,5}}
\end{tabproof}

\FormulaThmDeltaK[Die Multiplikation bildet ein Monoid]{%
  \MonoidAlg{\mathbb{N}}{\cdot}{1}%
}{NaturalMultiplicationMonoid}{%
  \DeltaRow{Mengen}{\mathbb{N}\dsep 1\dsep a}
}
\begin{tabproof}
  \proofstep{}{\SemigroupAlg{\mathbb{N}}{\cdot}}{%
    \FormulaRefAuto{NaturalMultiplicationSemigroup}[thm]}
  \proofstep{}{1\in\mathbb{N}}{%
    \FormulaRefAuto{PeanoOneInN}[thm]}
  \proofstep{3}{a\in\mathbb{N}}{\rA}
  \proofstep{3}{1\cdot a=a}{%
    \FormulaRefAuto{PeanoMulLeftOne}[thm]{3}}
  \proofstep{3}{a\cdot1=a}{%
    \FormulaRefAuto{PeanoMulRightOne}[thm]{3}}
  \proofstep{}{\MonoidAlg{\mathbb{N}}{\cdot}{1}}{%
    \FormulaRefAuto{MonoidDef}[def]{1,2,4,5}}
\end{tabproof}

\FormulaThmDeltaK[Die Multiplikation bildet ein kommutatives Monoid]{%
  \CommutativeMonoidAlg{\mathbb{N}}{\cdot}{1}%
}{NaturalMultiplicativeCommutativeMonoid}{%
  \DeltaRow{Mengen}{\mathbb{N}\dsep 1\dsep a\dsep b}
}
\begin{tabproof}
  \proofstep{}{\MonoidAlg{\mathbb{N}}{\cdot}{1}}{%
    \FormulaRefAuto{NaturalMultiplicationMonoid}[thm]}
  \proofstep{2}{a\in\mathbb{N}}{\rA}
  \proofstep{3}{b\in\mathbb{N}}{\rA}
  \proofstep{2,3}{a\cdot b=b\cdot a}{%
    \FormulaRefAuto{PeanoMulCommutative}[thm]{2,3}}
  \proofstep{}{\CommutativeMonoidAlg{\mathbb{N}}{\cdot}{1}}{%
    \FormulaRefAuto{CommutativeMonoidDef}[def]{1,4}}
\end{tabproof}

\section{Der kommutative Halbring mit Eins}

\FormulaThmDeltaK[Additive und multiplikative Monoidstruktur]{%
  \begin{aligned}
    &\CommutativeMonoidAlg{\mathbb{N}}{+}{0}\land{}\\[-2pt]
    &\MonoidAlg{\mathbb{N}}{\cdot}{1}
  \end{aligned}%
}{NaturalSemiringMonoidStructures}{%
  \DeltaRow{Mengen}{\mathbb{N}\dsep 0\dsep 1}
}
\begin{tabproofwide}
  \proofstepwidestar[]{%
    \CommutativeMonoidAlg{\mathbb{N}}{+}{0}}{%
    \FormulaRefAuto{NaturalAdditiveCommutativeMonoid}[thm]}
  \proofstepwidestar[]{%
    \CommutativeMonoidAlg{\mathbb{N}}{\cdot}{1}}{%
    \FormulaRefAuto{NaturalMultiplicativeCommutativeMonoid}[thm]}
  \proofstepwidestar[]{\MonoidAlg{\mathbb{N}}{\cdot}{1}}{%
    \FormulaRefAuto{CommutativeMonoidBaseMonoidAxiom}[ax]{2}}
  \proofstepwidestar[]{%
    \CommutativeMonoidAlg{\mathbb{N}}{+}{0}
    \land\MonoidAlg{\mathbb{N}}{\cdot}{1}}{%
    \rAI{1,3}}
\end{tabproofwide}

\FormulaThmDeltaK[Die natürlichen Zahlen bilden einen Halbring mit Eins]{%
  \SemiringAlg{\mathbb{N}}{+}{\cdot}{0}{1}%
}{NaturalUnitalSemiring}{%
  \DeltaRow{Mengen}{\mathbb{N}\dsep 0\dsep 1\dsep a\dsep b\dsep c}
}
\begin{tabproofwide}
  \proofstepwidestar[]{%
    \CommutativeMonoidAlg{\mathbb{N}}{+}{0}
    \land\MonoidAlg{\mathbb{N}}{\cdot}{1}}{%
    \FormulaRefAuto{NaturalSemiringMonoidStructures}[thm]}
  \proofstepwidestar[]{%
    \CommutativeMonoidAlg{\mathbb{N}}{+}{0}}{\rAEa{1}}
  \proofstepwidestar[]{%
    \MonoidAlg{\mathbb{N}}{\cdot}{1}}{\rAEb{1}}
  \proofstepwidestar[4]{a\in\mathbb{N}}{\rA}
  \proofstepwidestar[5]{b\in\mathbb{N}}{\rA}
  \proofstepwidestar[6]{c\in\mathbb{N}}{\rA}
  \proofstepwidestar[4,5,6]{%
    a\cdot(b+c)=a\cdot b+a\cdot c}{%
    \FormulaRefAuto{PeanoMulLeftDistributive}[thm]{4,5,6}}
  \proofstepwidestar[4,5,6]{%
    (a+b)\cdot c=a\cdot c+b\cdot c}{%
    \FormulaRefAuto{PeanoMulRightDistributive}[thm]{4,5,6}}
  \proofstepwidestar[4]{0\cdot a=0}{%
    \FormulaRefAuto{PeanoMulLeftZero}[thm]{4}}
  \proofstepwidestar[4]{a\cdot0=0}{%
    \FormulaRefAuto{PeanoMulRightZero}[thm]{4}}
  \proofstepwidestar[]{%
    \SemiringAlg{\mathbb{N}}{+}{\cdot}{0}{1}}{%
    \FormulaRefAuto{SemiringDef}[def]{2,3,7,8,9,10}}
\end{tabproofwide}

\FormulaThmDeltaK[Die natürlichen Zahlen bilden einen kommutativen Halbring mit Eins]{%
  \CommutativeSemiringAlg{\mathbb{N}}{+}{\cdot}{0}{1}%
}{NaturalCommutativeUnitalSemiring}{%
  \DeltaRow{Mengen}{\mathbb{N}\dsep 0\dsep 1}
}
\begin{tabproofwide}
  \proofstepwidestar[]{%
    \SemiringAlg{\mathbb{N}}{+}{\cdot}{0}{1}}{%
    \FormulaRefAuto{NaturalUnitalSemiring}[thm]}
  \proofstepwidestar[]{%
    \CommutativeMonoidAlg{\mathbb{N}}{\cdot}{1}}{%
    \FormulaRefAuto{NaturalMultiplicativeCommutativeMonoid}[thm]}
  \proofstepwidestar[]{%
    \CommutativeSemiringAlg{\mathbb{N}}{+}{\cdot}{0}{1}}{%
    \FormulaRefAuto{CommutativeUnitalSemiringDef}[def]{1,2}}
\end{tabproofwide}

\chapter{Die universelle Eigenschaft der natürlichen Zahlen}

\section{Die kanonische Abbildung}

\FormulaDefDeltaK[Nachfolgerabbildung eines Halbrings]{%
  \SemiringAlg{R}{\oplus}{\odot}{0_R}{1_R}\dsep x\in R
  \vdash
  \SemiringSuccessor{R,\oplus,\odot,0_R,1_R}(x)
  \coloneqq x\oplus1_R%
}{SemiringSuccessorDef}{%
  \DeltaRow{Mengen}{R\dsep0_R\dsep1_R\dsep x}
  \DeltaRow{Funktionen}{\oplus\dsep\odot}
  \DeltaRow{\textbf{Neues Symbol}}{}
  \DeltaPrem{Nachfolgerabbildung}{%
    \SemiringSuccessor{R,\oplus,\odot,0_R,1_R}}
}

\FormulaThmDeltaK[Die Halbringnachfolgerabbildung ist eine Endoabbildung]{%
  \SemiringAlg{R}{\oplus}{\odot}{0_R}{1_R}
  \vdash
  \SemiringSuccessor{R,\oplus,\odot,0_R,1_R}\colon R\to R%
}{SemiringSuccessorFunction}{%
  \DeltaRow{Mengen}{R\dsep0_R\dsep1_R\dsep x}
  \DeltaRow{Funktionen}{\oplus\dsep\odot}
}
\begin{tabproof}
  \proofstep{1}{\SemiringAlg{R}{\oplus}{\odot}{0_R}{1_R}}{\rA}
  \proofstep{1}{\CommutativeMonoidAlg{R}{\oplus}{0_R}}{%
    \FormulaRefAuto{SemiringAdditiveCommutativeMonoidAxiom}[ax]{1}}
  \proofstep{1}{\MonoidAlg{R}{\oplus}{0_R}}{%
    \FormulaRefAuto{CommutativeMonoidBaseMonoidAxiom}[ax]{2}}
  \proofstep{1}{\SemigroupAlg{R}{\oplus}}{%
    \FormulaRefAuto{MonoidBaseSemigroupAxiom}[ax]{3}}
  \proofstep{1}{\MonoidAlg{R}{\odot}{1_R}}{%
    \FormulaRefAuto{SemiringMultiplicativeMonoidAxiom}[ax]{1}}
  \proofstep{1}{1_R\in R}{%
    \FormulaRefAuto{MonoidIdentityCarrierAxiom}[ax]{5}}
  \proofstep{7}{x\in R}{\rA}
  \proofstep{1,7}{x\oplus1_R\in R}{%
    \FormulaRefAuto{SemigroupClosureAxiom}[ax]{4,7,6}}
  \proofstep{1,7}{%
    \SemiringSuccessor{R,\oplus,\odot,0_R,1_R}(x)=x\oplus1_R}{%
    \FormulaRefAuto{SemiringSuccessorDef}[def]{1,7}}
  \proofstep{1,7}{%
    \SemiringSuccessor{R,\oplus,\odot,0_R,1_R}(x)\in R}{%
    \rIE{9,8}}
  \proofstep{1}{%
    \SemiringSuccessor{R,\oplus,\odot,0_R,1_R}\colon R\to R}{%
    \FormulaRefAuto{%
      x\in A\vdash t(x)\coloneq y\dsep x\in A\vdash t(x)\in B
      \exists! F\colon A\to B\,\forall x\in A\,F(x)=t(x)}{9,10}}
\end{tabproof}

\FormulaThmDeltaK[Das additive Nullelement liegt im Träger]{%
  \SemiringAlg{R}{\oplus}{\odot}{0_R}{1_R}
  \vdash
  0_R\in R%
}{SemiringZeroCarrier}{%
  \DeltaRow{Mengen}{R\dsep0_R\dsep1_R}
  \DeltaRow{Funktionen}{\oplus\dsep\odot}
}
\begin{tabproof}
  \proofstep{1}{\SemiringAlg{R}{\oplus}{\odot}{0_R}{1_R}}{\rA}
  \proofstep{1}{\CommutativeMonoidAlg{R}{\oplus}{0_R}}{%
    \FormulaRefAuto{SemiringAdditiveCommutativeMonoidAxiom}[ax]{1}}
  \proofstep{1}{\MonoidAlg{R}{\oplus}{0_R}}{%
    \FormulaRefAuto{CommutativeMonoidBaseMonoidAxiom}[ax]{2}}
  \proofstep{1}{0_R\in R}{%
    \FormulaRefAuto{MonoidIdentityCarrierAxiom}[ax]{3}}
\end{tabproof}

\FormulaDefDeltaK[Kanonische Abbildung in einen Halbring]{%
  \begin{aligned}[t]
    &\SemiringAlg{R}{\oplus}{\odot}{0_R}{1_R}\vdash{}\\[-2pt]
    &\NaturalSemiringMap{R,\oplus,\odot,0_R,1_R}
      \coloneqq
      \RecFun{0_R}{\SemiringSuccessor{R,\oplus,\odot,0_R,1_R}}
  \end{aligned}
}{NaturalSemiringMapDef}{%
  \DeltaRow{Mengen}{R\dsep0_R\dsep1_R\dsep n}
  \DeltaRow{Funktionen}{\oplus\dsep\odot}
  \DeltaRow{Startwert}{}[%
    \FormulaRefAuto{SemiringZeroCarrier}[thm]]
  \DeltaRow{Schrittfunktion}{}[%
    \FormulaRefAuto{SemiringSuccessorFunction}[thm]]
  \DeltaRow{Rekursionssatz}{}[%
    \FormulaRefAuto{DedekindRecursionTheorem}[thm]]
  \DeltaRow{\textbf{Neues Symbol}}{}
  \DeltaPrem{Symbol}{\eta}
}

\FormulaThmDeltaK[Trägerabbildung der kanonischen Abbildung]{%
  \SemiringAlg{R}{\oplus}{\odot}{0_R}{1_R}
  \vdash
  \NaturalSemiringMap{R,\oplus,\odot,0_R,1_R}
    \colon\mathbb N\to R%
}{NaturalSemiringMapFunction}{%
  \DeltaRow{Mengen}{R\dsep0_R\dsep1_R}
  \DeltaRow{Funktionen}{\oplus\dsep\odot}
}
\begin{tabproof}
  \proofstep{1}{\SemiringAlg{R}{\oplus}{\odot}{0_R}{1_R}}{\rA}
  \proofstep{1}{0_R\in R}{%
    \FormulaRefAuto{SemiringZeroCarrier}[thm]{1}}
  \proofstep{1}{%
    \SemiringSuccessor{R,\oplus,\odot,0_R,1_R}\colon R\to R}{%
    \FormulaRefAuto{SemiringSuccessorFunction}[thm]{1}}
  \proofstep{1}{%
    \RecFun{0_R}{\SemiringSuccessor{R,\oplus,\odot,0_R,1_R}}
    \colon\mathbb N\to R}{%
    \FormulaRefAuto{%
      m_0\in M\dsep f\colon M\to M \vdash
      \RecFun{m_0}{f}\colon \mathbb{N}\to M}{2,3}}
  \proofstep{1}{%
    \NaturalSemiringMap{R,\oplus,\odot,0_R,1_R}
    =
    \RecFun{0_R}{\SemiringSuccessor{R,\oplus,\odot,0_R,1_R}}}{%
    \FormulaRefAuto{NaturalSemiringMapDef}[def]{1}}
  \proofstep{1}{%
    \NaturalSemiringMap{R,\oplus,\odot,0_R,1_R}
    \colon\mathbb N\to R}{\rIE{5,4}}
\end{tabproof}

\FormulaThmDeltaK[Nullwert der kanonischen Abbildung]{%
  \SemiringAlg{R}{\oplus}{\odot}{0_R}{1_R}
  \vdash
  \NaturalSemiringMap{R,\oplus,\odot,0_R,1_R}(0)=0_R%
}{NaturalSemiringMapZero}{%
  \DeltaRow{Mengen}{R\dsep0_R\dsep1_R}
  \DeltaRow{Funktionen}{\oplus\dsep\odot}
}
\begin{tabproof}
  \proofstep{1}{\SemiringAlg{R}{\oplus}{\odot}{0_R}{1_R}}{\rA}
  \proofstep{1}{0_R\in R}{%
    \FormulaRefAuto{SemiringZeroCarrier}[thm]{1}}
  \proofstep{1}{%
    \SemiringSuccessor{R,\oplus,\odot,0_R,1_R}\colon R\to R}{%
    \FormulaRefAuto{SemiringSuccessorFunction}[thm]{1}}
  \proofstep{1}{%
    \RecFun{0_R}{\SemiringSuccessor{R,\oplus,\odot,0_R,1_R}}(0)
    =0_R}{%
    \FormulaRefAuto{%
      m_0\in M\dsep f\colon M\to M \vdash
      \RecFun{m_0}{f}(0)=m_0}{2,3}}
  \proofstep{1}{%
    \NaturalSemiringMap{R,\oplus,\odot,0_R,1_R}
    =
    \RecFun{0_R}{\SemiringSuccessor{R,\oplus,\odot,0_R,1_R}}}{%
    \FormulaRefAuto{NaturalSemiringMapDef}[def]{1}}
  \proofstep{1}{%
    \NaturalSemiringMap{R,\oplus,\odot,0_R,1_R}(0)=0_R}{%
    \rIE{5,4}}
\end{tabproof}

\FormulaThmDeltaK[Nachfolgerwert der kanonischen Abbildung]{%
  \begin{aligned}[t]
    &\SemiringAlg{R}{\oplus}{\odot}{0_R}{1_R}\dsep n\in\mathbb N\\
    &\vdash
      \NaturalSemiringMap{R,\oplus,\odot,0_R,1_R}(\Succ(n))
      =
      \NaturalSemiringMap{R,\oplus,\odot,0_R,1_R}(n)\oplus1_R
  \end{aligned}
}{NaturalSemiringMapSuccessor}{%
  \DeltaRow{Mengen}{R\dsep0_R\dsep1_R\dsep n}
  \DeltaRow{Funktionen}{\oplus\dsep\odot}
}
\begin{notation*}
\[
  \eta_R:=\NaturalSemiringMap{R,\oplus,\odot,0_R,1_R},
  \qquad
  s_R:=\SemiringSuccessor{R,\oplus,\odot,0_R,1_R}.
\]
\end{notation*}
\begin{tabproofwide}
  \proofstepwidestar[1]{%
    \SemiringAlg{R}{\oplus}{\odot}{0_R}{1_R}}{\rA}
  \proofstepwidestar[2]{n\in\mathbb N}{\rA}
  \proofstepwidestar[1]{%
    0_R\in R}{%
    \FormulaRefAuto{SemiringZeroCarrier}[thm]{1}}
  \proofstepwidestar[1]{s_R\colon R\to R}{%
    \FormulaRefAuto{SemiringSuccessorFunction}[thm]{1}}
  \proofstepwidestar[1]{\eta_R=\RecFun{0_R}{s_R}}{%
    \FormulaRefAuto{NaturalSemiringMapDef}[def]{1}}
  \proofstepwidestar[1]{\eta_R\colon\mathbb N\to R}{%
    \FormulaRefAuto{NaturalSemiringMapFunction}[thm]{1}}
  \proofstepwidestar[1,2]{\eta_R(n)\in R}{%
    \FormulaRefAuto{F\colon A\to B\dsep a\in A\vdash F(a)\in B}{6,2}}
  \proofstepwidestar[1]{\RecFun{0_R}{s_R}=\eta_R}{%
    \FormulaRefAuto{a=b\vdash b=a}{5}}

  \proofstepwide[1,2]{\eta_R(\Succ(n))}{=}{%
    \RecFun{0_R}{s_R}(\Succ(n))}{\rIE{5}}
  \proofstepwide[1,2]{}{=}{%
    s_R\!\left(\RecFun{0_R}{s_R}(n)\right)}{%
    \FormulaRefAuto{%
      m_0\in M\dsep f\colon M\to M\dsep n\in\mathbb{N}\vdash
      \RecFun{m_0}{f}(\Succ(n))=f(\RecFun{m_0}{f}(n))}{3,4,2}}
  \proofstepwide[1,2]{}{=}{s_R\!\left(\eta_R(n)\right)}{\rIE{8}}
  \proofstepwide[1,2]{}{=}{\eta_R(n)\oplus1_R}{%
    \FormulaRefAuto{SemiringSuccessorDef}[def]{1,7}}
  \proofstepwidestar[1,2]{%
    \eta_R(\Succ(n))=\eta_R(n)\oplus1_R}{%
    \rChain{9,12}}
\end{tabproofwide}

\FormulaThmDeltaK[Einserhaltung der kanonischen Abbildung]{%
  \SemiringAlg{R}{\oplus}{\odot}{0_R}{1_R}
  \vdash
  \NaturalSemiringMap{R,\oplus,\odot,0_R,1_R}(1)=1_R%
}{NaturalSemiringMapOne}{%
  \DeltaRow{Mengen}{R\dsep0_R\dsep1_R}
  \DeltaRow{Funktionen}{\oplus\dsep\odot}
}
\begin{tabproofwide}
  \proofstepwidestar[1]{%
    \SemiringAlg{R}{\oplus}{\odot}{0_R}{1_R}}{\rA}
  \proofstepwidestar[]{0\in\mathbb N}{\FormulaRefAuto{PeanoZero}[ax]}
  \proofstepwidestar[]{1=\Succ(0)}{\FormulaRefAuto{PeanoOne}[def]}
  \proofstepwidestar[1]{%
    \NaturalSemiringMap{R,\oplus,\odot,0_R,1_R}(0)=0_R}{%
    \FormulaRefAuto{NaturalSemiringMapZero}[thm]{1}}
  \proofstepwidestar[1]{\CommutativeMonoidAlg{R}{\oplus}{0_R}}{%
    \FormulaRefAuto{SemiringAdditiveCommutativeMonoidAxiom}[ax]{1}}
  \proofstepwidestar[1]{\MonoidAlg{R}{\oplus}{0_R}}{%
    \FormulaRefAuto{CommutativeMonoidBaseMonoidAxiom}[ax]{5}}
  \proofstepwidestar[1]{\MonoidAlg{R}{\odot}{1_R}}{%
    \FormulaRefAuto{SemiringMultiplicativeMonoidAxiom}[ax]{1}}
  \proofstepwidestar[1]{1_R\in R}{%
    \FormulaRefAuto{MonoidIdentityCarrierAxiom}[ax]{7}}

  \proofstepwide[1]{%
    \NaturalSemiringMap{R,\oplus,\odot,0_R,1_R}(1)}{=}{%
    \NaturalSemiringMap{R,\oplus,\odot,0_R,1_R}(\Succ(0))}{\rIE{3}}
  \proofstepwide[1]{}{=}{%
    \NaturalSemiringMap{R,\oplus,\odot,0_R,1_R}(0)\oplus1_R}{%
    \FormulaRefAuto{NaturalSemiringMapSuccessor}[thm]{1,2}}
  \proofstepwide[1]{}{=}{0_R\oplus1_R}{\rIE{4}}
  \proofstepwide[1]{}{=}{1_R}{%
    \FormulaRefAuto{MonoidLeftIdentityAxiom}[ax]{6,8}}
  \proofstepwidestar[1]{%
    \NaturalSemiringMap{R,\oplus,\odot,0_R,1_R}(1)=1_R}{%
    \rChain{9,12}}
\end{tabproofwide}

\begin{notation*}
  \[
    \eta_R\coloneqq
    \NaturalSemiringMap{R,\oplus,\odot,0_R,1_R}.
  \]
\end{notation*}

\FormulaThmDeltaK[Nullfall der Additionserhaltung]{%
  \begin{aligned}[t]
    &\SemiringAlg{R}{\oplus}{\odot}{0_R}{1_R}\dsep m\in\mathbb N\\
    &\vdash
      \NaturalSemiringMap{R,\oplus,\odot,0_R,1_R}(m+0)
      =
      \NaturalSemiringMap{R,\oplus,\odot,0_R,1_R}(m)
      \oplus
      \NaturalSemiringMap{R,\oplus,\odot,0_R,1_R}(0)
  \end{aligned}
}{NaturalSemiringMapAdditionZero}{%
  \DeltaRow{Mengen}{R\dsep0_R\dsep1_R\dsep m}
  \DeltaRow{Funktionen}{\oplus\dsep\odot}
}
\begin{tabproofwide}
  \proofstepwidestar[1]{%
    \SemiringAlg{R}{\oplus}{\odot}{0_R}{1_R}}{\rA}
  \proofstepwidestar[2]{m\in\mathbb N}{\rA}
  \proofstepwidestar[1]{%
    \CommutativeMonoidAlg{R}{\oplus}{0_R}}{%
    \FormulaRefAuto{SemiringAdditiveCommutativeMonoidAxiom}[ax]{1}}
  \proofstepwidestar[1]{\MonoidAlg{R}{\oplus}{0_R}}{%
    \FormulaRefAuto{CommutativeMonoidBaseMonoidAxiom}[ax]{3}}
  \proofstepwidestar[1]{\eta_R\colon\mathbb N\to R}{%
    \FormulaRefAuto{NaturalSemiringMapFunction}[thm]{1}}
  \proofstepwidestar[1,2]{\eta_R(m)\in R}{%
    \FormulaRefAuto{F\colon A\to B\dsep a\in A\vdash F(a)\in B}{5,2}}
  \proofstepwide[1,2]{\eta_R(m)\oplus0_R}{=}{\eta_R(m)}{%
    \FormulaRefAuto{MonoidRightIdentityAxiom}[ax]{4,6}}
  \proofstepwidestar[1]{\eta_R(0)=0_R}{%
    \FormulaRefAuto{NaturalSemiringMapZero}[thm]{1}}
  \proofstepwidestar[1]{0_R=\eta_R(0)}{%
    \FormulaRefAuto{a=b\vdash b=a}{8}}
  \proofstepwide[2]{\eta_R(m+0)}{=}{\eta_R(m)}{%
    \FormulaRefAuto{PeanoAddRightZero}[thm]{2}}
  \proofstepwide[1,2]{}{=}{\eta_R(m)\oplus0_R}{%
    \FormulaRefAuto{a=b\vdash b=a}{7}}
  \proofstepwide[1,2]{}{=}{\eta_R(m)\oplus\eta_R(0)}{\rIE{9}}
  \proofstepwidestar[1,2]{%
    \eta_R(m+0)=\eta_R(m)\oplus\eta_R(0)}{%
    \rChain{10,12}}
\end{tabproofwide}

\FormulaThmDeltaK[Nachfolgerschritt der Additionserhaltung]{%
  \begin{aligned}[t]
    &\SemiringAlg{R}{\oplus}{\odot}{0_R}{1_R}\dsep m,n\in\mathbb N\dsep{}\\[-2pt]
    &\NaturalSemiringMap{R,\oplus,\odot,0_R,1_R}(m+n)
      =
      \NaturalSemiringMap{R,\oplus,\odot,0_R,1_R}(m)
      \oplus
      \NaturalSemiringMap{R,\oplus,\odot,0_R,1_R}(n)\\
    &\vdash
      \NaturalSemiringMap{R,\oplus,\odot,0_R,1_R}(m+\Succ(n))
      =
      \NaturalSemiringMap{R,\oplus,\odot,0_R,1_R}(m)
      \oplus
      \NaturalSemiringMap{R,\oplus,\odot,0_R,1_R}(\Succ(n))
  \end{aligned}
}{NaturalSemiringMapAdditionStep}{%
  \DeltaRow{Mengen}{R\dsep0_R\dsep1_R\dsep m\dsep n}
  \DeltaRow{Funktionen}{\oplus\dsep\odot}
}
\begin{tabproofwide}
  \proofstepwidestar[1]{%
    \SemiringAlg{R}{\oplus}{\odot}{0_R}{1_R}}{\rA}
  \proofstepwidestar[2]{m\in\mathbb N}{\rA}
  \proofstepwidestar[3]{n\in\mathbb N}{\rA}
  \proofstepwidestar[4]{\eta_R(m+n)=\eta_R(m)\oplus\eta_R(n)}{\rA}
  \proofstepwidestar[1]{%
    \CommutativeMonoidAlg{R}{\oplus}{0_R}}{%
    \FormulaRefAuto{SemiringAdditiveCommutativeMonoidAxiom}[ax]{1}}
  \proofstepwidestar[1]{\MonoidAlg{R}{\oplus}{0_R}}{%
    \FormulaRefAuto{CommutativeMonoidBaseMonoidAxiom}[ax]{5}}
  \proofstepwidestar[1]{\SemigroupAlg{R}{\oplus}}{%
    \FormulaRefAuto{MonoidBaseSemigroupAxiom}[ax]{6}}
  \proofstepwidestar[1]{\eta_R\colon\mathbb N\to R}{%
    \FormulaRefAuto{NaturalSemiringMapFunction}[thm]{1}}
  \proofstepwidestar[2,3]{m+n\in\mathbb N}{%
    \FormulaRefAuto{PeanoAddClosure}[thm]{2,3}}
  \proofstepwidestar[1,2]{\eta_R(m)\in R}{%
    \FormulaRefAuto{F\colon A\to B\dsep a\in A\vdash F(a)\in B}{8,2}}
  \proofstepwidestar[1,3]{\eta_R(n)\in R}{%
    \FormulaRefAuto{F\colon A\to B\dsep a\in A\vdash F(a)\in B}{8,3}}
  \proofstepwidestar[1]{\MonoidAlg{R}{\odot}{1_R}}{%
    \FormulaRefAuto{SemiringMultiplicativeMonoidAxiom}[ax]{1}}
  \proofstepwidestar[1]{1_R\in R}{%
    \FormulaRefAuto{MonoidIdentityCarrierAxiom}[ax]{12}}
  \proofstepwidestar[1,3]{\eta_R(\Succ(n))=\eta_R(n)\oplus1_R}{%
    \FormulaRefAuto{NaturalSemiringMapSuccessor}[thm]{1,3}}
  \proofstepwidestar[1,3]{\eta_R(n)\oplus1_R=\eta_R(\Succ(n))}{%
    \FormulaRefAuto{a=b\vdash b=a}{14}}
  \proofstepwide[2,3]{\eta_R(m+\Succ(n))}{=}{\eta_R(\Succ(m+n))}{%
    \FormulaRefAuto{PeanoAddSuccRight}[thm]{2,3}}
  \proofstepwide[1,2,3]{}{=}{\eta_R(m+n)\oplus1_R}{%
    \FormulaRefAuto{NaturalSemiringMapSuccessor}[thm]{1,9}}
  \proofstepwide[1,2,3,4]{}{=}{%
    \bigl(\eta_R(m)\oplus\eta_R(n)\bigr)\oplus1_R}{\rIE{4}}
  \proofstepwide[1,2,3]{}{=}{\eta_R(m)\oplus\bigl(\eta_R(n)\oplus1_R\bigr)}{%
    \FormulaRefAuto{SemigroupAssociativityAxiom}[ax]{7,10,11,13}}
  \proofstepwide[1,2,3]{}{=}{\eta_R(m)\oplus\eta_R(\Succ(n))}{\rIE{15}}
  \proofstepwidestar[1,2,3,4]{%
    \eta_R(m+\Succ(n))=\eta_R(m)\oplus\eta_R(\Succ(n))}{%
    \rChain{16,20}}
\end{tabproofwide}

\FormulaThmDeltaK[Die kanonische Abbildung erhält die Addition]{%
  \begin{aligned}[t]
    &\SemiringAlg{R}{\oplus}{\odot}{0_R}{1_R}\dsep m,n\in\mathbb N\\
    &\vdash
      \NaturalSemiringMap{R,\oplus,\odot,0_R,1_R}(m+n)
      =
      \NaturalSemiringMap{R,\oplus,\odot,0_R,1_R}(m)
      \oplus
      \NaturalSemiringMap{R,\oplus,\odot,0_R,1_R}(n)
  \end{aligned}
}{NaturalSemiringMapAddition}{%
  \DeltaRow{Mengen}{R\dsep0_R\dsep1_R\dsep m\dsep n\dsep k}
  \DeltaRow{Funktionen}{\oplus\dsep\odot}
}
\begin{tabproofwide}
  \proofstepwidestar[1]{%
    \SemiringAlg{R}{\oplus}{\odot}{0_R}{1_R}}{\rA}
  \proofstepwidestar[2]{m\in\mathbb N}{\rA}
  \proofstepwidestar[3]{n\in\mathbb N}{\rA}
  \proofstepwidestar[1,2]{\eta_R(m+0)=\eta_R(m)\oplus\eta_R(0)}{%
    \FormulaRefAuto{NaturalSemiringMapAdditionZero}[thm]{1,2}}
  \proofstepwidestar[5]{k\in\mathbb N}{\rA}
  \proofstepwidestar[6]{\eta_R(m+k)=\eta_R(m)\oplus\eta_R(k)}{\rA}
  \proofstepwidestar[1,2,5,6]{%
    \eta_R(m+\Succ(k))=\eta_R(m)\oplus\eta_R(\Succ(k))}{%
    \FormulaRefAuto{NaturalSemiringMapAdditionStep}[thm]{1,2,5,6}}
  \proofstepwidestar[1,2,3]{\eta_R(m+n)=\eta_R(m)\oplus\eta_R(n)}{%
    \FormulaRefAuto{PeanoInductionRule}[ax]{4,7,3}}
\end{tabproofwide}

\FormulaThmDeltaK[Nullfall der Multiplikationserhaltung]{%
  \begin{aligned}[t]
    &\SemiringAlg{R}{\oplus}{\odot}{0_R}{1_R}\dsep m\in\mathbb N\\
    &\vdash
      \NaturalSemiringMap{R,\oplus,\odot,0_R,1_R}(m\cdot0)
      =
      \NaturalSemiringMap{R,\oplus,\odot,0_R,1_R}(m)
      \odot
      \NaturalSemiringMap{R,\oplus,\odot,0_R,1_R}(0)
  \end{aligned}
}{NaturalSemiringMapMultiplicationZero}{%
  \DeltaRow{Mengen}{R\dsep0_R\dsep1_R\dsep m}
  \DeltaRow{Funktionen}{\oplus\dsep\odot}
}
\begin{tabproofwide}
  \proofstepwidestar[1]{%
    \SemiringAlg{R}{\oplus}{\odot}{0_R}{1_R}}{\rA}
  \proofstepwidestar[2]{m\in\mathbb N}{\rA}
  \proofstepwidestar[1]{\eta_R\colon\mathbb N\to R}{%
    \FormulaRefAuto{NaturalSemiringMapFunction}[thm]{1}}
  \proofstepwidestar[1,2]{\eta_R(m)\in R}{%
    \FormulaRefAuto{F\colon A\to B\dsep a\in A\vdash F(a)\in B}{3,2}}
  \proofstepwidestar[1,2]{\eta_R(m)\odot0_R=0_R}{%
    \FormulaRefAuto{SemiringRightZeroAxiom}[ax]{1,4}}
  \proofstepwidestar[1,2]{0_R=\eta_R(m)\odot0_R}{%
    \FormulaRefAuto{a=b\vdash b=a}{5}}
  \proofstepwidestar[1]{\eta_R(0)=0_R}{%
    \FormulaRefAuto{NaturalSemiringMapZero}[thm]{1}}
  \proofstepwidestar[1]{0_R=\eta_R(0)}{%
    \FormulaRefAuto{a=b\vdash b=a}{7}}
  \proofstepwide[2]{\eta_R(m\cdot0)}{=}{\eta_R(0)}{%
    \FormulaRefAuto{PeanoMulRightZero}[thm]{2}}
  \proofstepwide[1,2]{}{=}{0_R}{\rIE{7}}
  \proofstepwide[1,2]{}{=}{\eta_R(m)\odot0_R}{\rIE{6}}
  \proofstepwide[1,2]{}{=}{\eta_R(m)\odot\eta_R(0)}{\rIE{8}}
  \proofstepwidestar[1,2]{%
    \eta_R(m\cdot0)=\eta_R(m)\odot\eta_R(0)}{%
    \rChain{9,12}}
\end{tabproofwide}

\FormulaThmDeltaK[Nachfolgerschritt der Multiplikationserhaltung]{%
  \begin{aligned}[t]
    &\SemiringAlg{R}{\oplus}{\odot}{0_R}{1_R}\dsep m,n\in\mathbb N\dsep{}\\[-2pt]
    &\NaturalSemiringMap{R,\oplus,\odot,0_R,1_R}(m\cdot n)
      =
      \NaturalSemiringMap{R,\oplus,\odot,0_R,1_R}(m)
      \odot
      \NaturalSemiringMap{R,\oplus,\odot,0_R,1_R}(n)\\
    &\vdash
      \NaturalSemiringMap{R,\oplus,\odot,0_R,1_R}(m\cdot\Succ(n))
      =
      \NaturalSemiringMap{R,\oplus,\odot,0_R,1_R}(m)
      \odot
      \NaturalSemiringMap{R,\oplus,\odot,0_R,1_R}(\Succ(n))
  \end{aligned}
}{NaturalSemiringMapMultiplicationStep}{%
  \DeltaRow{Mengen}{R\dsep0_R\dsep1_R\dsep m\dsep n}
  \DeltaRow{Funktionen}{\oplus\dsep\odot}
}
\begin{notation*}
\[
  x:=\eta_R(m),
  \qquad
  y:=\eta_R(n).
\]
\end{notation*}
\begin{tabproofwide}
  \proofstepwidestar[1]{%
    \SemiringAlg{R}{\oplus}{\odot}{0_R}{1_R}}{\rA}
  \proofstepwidestar[2]{m\in\mathbb N}{\rA}
  \proofstepwidestar[3]{n\in\mathbb N}{\rA}
  \proofstepwidestar[4]{\eta_R(m\cdot n)=x\odot y}{\rA}
  \proofstepwidestar[1]{\eta_R\colon\mathbb N\to R}{%
    \FormulaRefAuto{NaturalSemiringMapFunction}[thm]{1}}
  \proofstepwidestar[2,3]{m\cdot n\in\mathbb N}{%
    \FormulaRefAuto{PeanoMulClosure}[thm]{2,3}}
  \proofstepwidestar[1,2]{x\in R}{%
    \FormulaRefAuto{F\colon A\to B\dsep a\in A\vdash F(a)\in B}{5,2}}
  \proofstepwidestar[1,3]{y\in R}{%
    \FormulaRefAuto{F\colon A\to B\dsep a\in A\vdash F(a)\in B}{5,3}}
  \proofstepwidestar[1]{\MonoidAlg{R}{\odot}{1_R}}{%
    \FormulaRefAuto{SemiringMultiplicativeMonoidAxiom}[ax]{1}}
  \proofstepwidestar[1]{1_R\in R}{%
    \FormulaRefAuto{MonoidIdentityCarrierAxiom}[ax]{9}}
  \proofstepwidestar[1,2]{x\odot1_R=x}{%
    \FormulaRefAuto{MonoidRightIdentityAxiom}[ax]{9,7}}
  \proofstepwidestar[1,2]{x=x\odot1_R}{%
    \FormulaRefAuto{a=b\vdash b=a}{11}}
  \proofstepwidestar[1,2,3]{%
    \begin{aligned}[t]
      &x\odot\bigl(y\oplus1_R\bigr)\\[-2pt]
      &\quad=
        \bigl(x\odot y\bigr)\oplus\bigl(x\odot1_R\bigr)
    \end{aligned}}{%
    \FormulaRefAuto{SemiringLeftDistributiveAxiom}[ax]{1,7,8,10}}
  \proofstepwidestar[1,2,3]{%
    \begin{aligned}[t]
      &\bigl(x\odot y\bigr)\oplus\bigl(x\odot1_R\bigr)\\[-2pt]
      &\quad=x\odot\bigl(y\oplus1_R\bigr)
    \end{aligned}}{%
    \FormulaRefAuto{a=b\vdash b=a}{13}}
  \proofstepwidestar[1,3]{\eta_R(\Succ(n))=y\oplus1_R}{%
    \FormulaRefAuto{NaturalSemiringMapSuccessor}[thm]{1,3}}
  \proofstepwidestar[1,3]{y\oplus1_R=\eta_R(\Succ(n))}{%
    \FormulaRefAuto{a=b\vdash b=a}{15}}
  \proofstepwide[2,3]{\eta_R(m\cdot\Succ(n))}{=}{\eta_R(m\cdot n+m)}{%
    \FormulaRefAuto{PeanoMulSuccRight}[thm]{2,3}}
  \proofstepwide[1,2,3]{}{=}{\eta_R(m\cdot n)\oplus x}{%
    \FormulaRefAuto{NaturalSemiringMapAddition}[thm]{1,6,2}}
  \proofstepwide[1,2,3,4]{}{=}{%
    \bigl(x\odot y\bigr)\oplus x}{\rIE{4}}
  \proofstepwide[1,2,3,4]{}{=}{%
    \bigl(x\odot y\bigr)\oplus\bigl(x\odot1_R\bigr)}{\rIE{12}}
  \proofstepwide[1,2,3,4]{}{=}{%
    x\odot\bigl(y\oplus1_R\bigr)}{\rIE{14}}
  \proofstepwide[1,2,3,4]{}{=}{x\odot\eta_R(\Succ(n))}{\rIE{16}}
  \proofstepwidestar[1,2,3,4]{%
    \begin{aligned}[t]
      &\eta_R(m\cdot\Succ(n))\\[-2pt]
      &\quad=x\odot\eta_R(\Succ(n))
    \end{aligned}}{%
    \rChain{17,22}}
\end{tabproofwide}

\FormulaThmDeltaK[Die kanonische Abbildung erhält die Multiplikation]{%
  \begin{aligned}[t]
    &\SemiringAlg{R}{\oplus}{\odot}{0_R}{1_R}\dsep m,n\in\mathbb N\\
    &\vdash
      \NaturalSemiringMap{R,\oplus,\odot,0_R,1_R}(m\cdot n)
      =
      \NaturalSemiringMap{R,\oplus,\odot,0_R,1_R}(m)
      \odot
      \NaturalSemiringMap{R,\oplus,\odot,0_R,1_R}(n)
  \end{aligned}
}{NaturalSemiringMapMultiplication}{%
  \DeltaRow{Mengen}{R\dsep0_R\dsep1_R\dsep k\dsep m\dsep n}
  \DeltaRow{Funktionen}{\oplus\dsep\odot}
}
\begin{tabproofwide}
  \proofstepwidestar[1]{%
    \SemiringAlg{R}{\oplus}{\odot}{0_R}{1_R}}{\rA}
  \proofstepwidestar[2]{m\in\mathbb N}{\rA}
  \proofstepwidestar[3]{n\in\mathbb N}{\rA}
  \proofstepwidestar[1,2]{\eta_R(m\cdot0)=\eta_R(m)\odot\eta_R(0)}{%
    \FormulaRefAuto{NaturalSemiringMapMultiplicationZero}[thm]{1,2}}
  \proofstepwidestar[5]{k\in\mathbb N}{\rA}
  \proofstepwidestar[6]{\eta_R(m\cdot k)=\eta_R(m)\odot\eta_R(k)}{\rA}
  \proofstepwidestar[1,2,5,6]{%
    \eta_R(m\cdot\Succ(k))=\eta_R(m)\odot\eta_R(\Succ(k))}{%
    \FormulaRefAuto{NaturalSemiringMapMultiplicationStep}[thm]{1,2,5,6}}
  \proofstepwidestar[1,2,3]{\eta_R(m\cdot n)=\eta_R(m)\odot\eta_R(n)}{%
    \FormulaRefAuto{PeanoInductionRule}[ax]{4,7,3}}
\end{tabproofwide}

\begin{notation*}
  \[
    \eta_R\coloneqq
    \NaturalSemiringMap{R,\oplus,\odot,0_R,1_R}.
  \]
\end{notation*}

\FormulaThmDeltaK[Additiver Monoidanteil der kanonischen Abbildung]{%
  \begin{aligned}[t]
    &\SemiringAlg{R}{\oplus}{\odot}{0_R}{1_R}\vdash{}\\[-2pt]
    &\MonoidHom{\NaturalSemiringMap{R,\oplus,\odot,0_R,1_R}}
      {\mathbb N,+,0}{R,\oplus,0_R}
  \end{aligned}
}{NaturalSemiringCanonicalAdditiveMonoidHom}{%
  \DeltaRow{Mengen}{R\dsep0_R\dsep1_R\dsep m\dsep n}
  \DeltaRow{Funktionen}{\oplus\dsep\odot}
}
\begin{tabproofwide}
    \proofstepwidestar[1]{%
      \SemiringAlg{R}{\oplus}{\odot}{0_R}{1_R}}{\rA}
    \proofstepwidestar[]{\MonoidAlg{\mathbb N}{+}{0}}{%
      \FormulaRefAuto{NaturalAdditionMonoid}[thm]}
    \proofstepwidestar[1]{\MonoidAlg{R}{\oplus}{0_R}}{%
      \FormulaRefAuto{CommutativeMonoidBaseMonoidAxiom}[ax]{%
        \FormulaRefAuto{SemiringAdditiveCommutativeMonoidAxiom}[ax]{1}}}
    \proofstepwidestar[]{\SemigroupAlg{\mathbb N}{+}}{%
      \FormulaRefAuto{MonoidBaseSemigroupAxiom}[ax]{2}}
    \proofstepwidestar[1]{\SemigroupAlg{R}{\oplus}}{%
      \FormulaRefAuto{MonoidBaseSemigroupAxiom}[ax]{3}}
    \proofstepwidestar[1]{\eta_R\colon\mathbb N\to R}{%
      \FormulaRefAuto{NaturalSemiringMapFunction}[thm]{1}}
    \proofstepwidestar[7]{m\in\mathbb N}{\rA}
    \proofstepwidestar[8]{n\in\mathbb N}{\rA}
    \proofstepwidestar[1,7,8]{%
      \eta_R(m+n)=\eta_R(m)\oplus\eta_R(n)}{%
      \FormulaRefAuto{NaturalSemiringMapAddition}[thm]{1,7,8}}
    \proofstepwidestar[1]{%
      \SemigroupHom{\eta_R}{\mathbb N,+}{R,\oplus}}{%
      \FormulaRefAuto{SemigroupHomDef}[def]{4,5,6,9}}
    \proofstepwidestar[1]{\eta_R(0)=0_R}{%
      \FormulaRefAuto{NaturalSemiringMapZero}[thm]{1}}
    \proofstepwidestar[1]{%
      \MonoidHom{\eta_R}{\mathbb N,+,0}{R,\oplus,0_R}}{%
      \FormulaRefAuto{MonoidHomDef}[def]{2,3,10,11}}
\end{tabproofwide}

\FormulaThmDeltaK[Multiplikativer Monoidanteil der kanonischen Abbildung]{%
  \begin{aligned}[t]
    &\SemiringAlg{R}{\oplus}{\odot}{0_R}{1_R}\vdash{}\\[-2pt]
    &\MonoidHom{\NaturalSemiringMap{R,\oplus,\odot,0_R,1_R}}
      {\mathbb N,\cdot,1}{R,\odot,1_R}
  \end{aligned}
}{NaturalSemiringCanonicalMultiplicativeMonoidHom}{%
  \DeltaRow{Mengen}{R\dsep0_R\dsep1_R\dsep m\dsep n}
  \DeltaRow{Funktionen}{\oplus\dsep\odot}
}
\begin{tabproofwide}
    \proofstepwidestar[1]{%
      \SemiringAlg{R}{\oplus}{\odot}{0_R}{1_R}}{\rA}
    \proofstepwidestar[]{\MonoidAlg{\mathbb N}{\cdot}{1}}{%
      \FormulaRefAuto{NaturalMultiplicationMonoid}[thm]}
    \proofstepwidestar[1]{\MonoidAlg{R}{\odot}{1_R}}{%
      \FormulaRefAuto{SemiringMultiplicativeMonoidAxiom}[ax]{1}}
    \proofstepwidestar[]{\SemigroupAlg{\mathbb N}{\cdot}}{%
      \FormulaRefAuto{MonoidBaseSemigroupAxiom}[ax]{2}}
    \proofstepwidestar[1]{\SemigroupAlg{R}{\odot}}{%
      \FormulaRefAuto{MonoidBaseSemigroupAxiom}[ax]{3}}
    \proofstepwidestar[1]{\eta_R\colon\mathbb N\to R}{%
      \FormulaRefAuto{NaturalSemiringMapFunction}[thm]{1}}
    \proofstepwidestar[7]{m\in\mathbb N}{\rA}
    \proofstepwidestar[8]{n\in\mathbb N}{\rA}
    \proofstepwidestar[1,7,8]{%
      \eta_R(m\cdot n)=\eta_R(m)\odot\eta_R(n)}{%
      \FormulaRefAuto{NaturalSemiringMapMultiplication}[thm]{1,7,8}}
    \proofstepwidestar[1]{%
      \SemigroupHom{\eta_R}{\mathbb N,\cdot}{R,\odot}}{%
      \FormulaRefAuto{SemigroupHomDef}[def]{4,5,6,9}}
    \proofstepwidestar[1]{\eta_R(1)=1_R}{%
      \FormulaRefAuto{NaturalSemiringMapOne}[thm]{1}}
    \proofstepwidestar[1]{%
      \MonoidHom{\eta_R}{\mathbb N,\cdot,1}{R,\odot,1_R}}{%
      \FormulaRefAuto{MonoidHomDef}[def]{2,3,10,11}}
\end{tabproofwide}

\FormulaThmDeltaK[Die kanonische Abbildung ist ein Halbringhomomorphismus]{%
  \SemiringAlg{R}{\oplus}{\odot}{0_R}{1_R}
  \vdash
  \SemiringHom{\NaturalSemiringMap{R,\oplus,\odot,0_R,1_R}}
    {\mathbb N,+,\cdot,0,1}{R,\oplus,\odot,0_R,1_R}%
}{NaturalSemiringCanonicalHomomorphism}{%
  \DeltaRow{Mengen}{R\dsep0_R\dsep1_R}
  \DeltaRow{Funktionen}{\oplus\dsep\odot}
}
\begin{tabproofwide}
    \proofstepwidestar[1]{%
      \SemiringAlg{R}{\oplus}{\odot}{0_R}{1_R}}{\rA}
    \proofstepwidestar[]{%
      \SemiringAlg{\mathbb N}{+}{\cdot}{0}{1}}{%
      \FormulaRefAuto{NaturalUnitalSemiring}[thm]}
    \proofstepwidestar[1]{%
      \MonoidHom{\NaturalSemiringMap{R,\oplus,\odot,0_R,1_R}}
        {\mathbb N,+,0}{R,\oplus,0_R}}{%
      \FormulaRefAuto{NaturalSemiringCanonicalAdditiveMonoidHom}[thm]{1}}
    \proofstepwidestar[1]{%
      \MonoidHom{\NaturalSemiringMap{R,\oplus,\odot,0_R,1_R}}
        {\mathbb N,\cdot,1}{R,\odot,1_R}}{%
      \FormulaRefAuto{NaturalSemiringCanonicalMultiplicativeMonoidHom}[thm]{1}}
    \proofstepwidestar[1]{%
      \SemiringHom{\NaturalSemiringMap{R,\oplus,\odot,0_R,1_R}}
        {\mathbb N,+,\cdot,0,1}{R,\oplus,\odot,0_R,1_R}}{%
      \FormulaRefAuto{SemiringHomDef}[def]{2,1,3,4}}
\end{tabproofwide}

\FormulaThmDeltaK[Eindeutigkeit auf jedem natürlichen Argument]{%
  \begin{aligned}[t]
    &\SemiringAlg{R}{\oplus}{\odot}{0_R}{1_R}\dsep
      \SemiringHom{f}{\mathbb N,+,\cdot,0,1}
        {R,\oplus,\odot,0_R,1_R}\dsep n\in\mathbb N\\
    &\vdash
      f(n)=\NaturalSemiringMap{R,\oplus,\odot,0_R,1_R}(n)
  \end{aligned}
}{NaturalSemiringHomPointwiseUnique}{%
  \DeltaRow{Mengen}{R\dsep0_R\dsep1_R\dsep n}
  \DeltaRow{Funktionen}{f\dsep\oplus\dsep\odot}
}
\begin{tabproofwide}
  \proofstepwidestar[1]{%
    \SemiringAlg{R}{\oplus}{\odot}{0_R}{1_R}}{\rA}
  \proofstepwidestar[2]{%
    \SemiringHom{f}{\mathbb N,+,\cdot,0,1}
      {R,\oplus,\odot,0_R,1_R}}{\rA}
  \proofstepwidestar[2]{f(0)=0_R\land f(1)=1_R}{%
    \FormulaRefAuto{SemiringHomConstantsAxiom}[ax]{2}}
  \proofstepwidestar[1]{\eta_R(0)=0_R}{%
    \FormulaRefAuto{NaturalSemiringMapZero}[thm]{1}}
  \proofstepwidestar[1]{0_R=\eta_R(0)}{%
    \FormulaRefAuto{a=b\vdash b=a}{4}}
  \proofstepwide[2]{f(0)}{=}{0_R}{\rAEa{3}}
  \proofstepwide[1,2]{}{=}{\eta_R(0)}{\rIE{5}}
  \proofstepwidestar[1,2]{f(0)=\eta_R(0)}{\rChain{6,7}}
  \proofstepwidestar[9]{n\in\mathbb N}{\rA}
  \proofstepwidestar[10]{f(n)=\eta_R(n)}{\rA}
  \proofstepwidestar[9]{n+1=\Succ(n)}{%
    \FormulaRefAuto{PeanoAddOneSucc}[thm]{9}}
  \proofstepwidestar[9]{\Succ(n)=n+1}{%
    \FormulaRefAuto{a=b\vdash b=a}{11}}
  \proofstepwidestar[]{1\in\mathbb N}{%
    \FormulaRefAuto{PeanoOneInN}[thm]}
  \proofstepwidestar[2]{f(1)=1_R}{\rAEb{3}}
  \proofstepwidestar[1,9]{\eta_R(\Succ(n))=\eta_R(n)\oplus1_R}{%
    \FormulaRefAuto{NaturalSemiringMapSuccessor}[thm]{1,9}}
  \proofstepwidestar[1,9]{\eta_R(n)\oplus1_R=\eta_R(\Succ(n))}{%
    \FormulaRefAuto{a=b\vdash b=a}{15}}
  \proofstepwide[9]{f(\Succ(n))}{=}{f(n+1)}{\rIE{12}}
  \proofstepwide[2,9]{}{=}{f(n)\oplus f(1)}{%
    \FormulaRefAuto{SemiringHomAdditionAxiom}[ax]{2,9,13}}
  \proofstepwide[1,2,9,10]{}{=}{\eta_R(n)\oplus f(1)}{\rIE{10}}
  \proofstepwide[1,2,9,10]{}{=}{\eta_R(n)\oplus1_R}{\rIE{14}}
  \proofstepwide[1,2,9,10]{}{=}{\eta_R(\Succ(n))}{\rIE{16}}
  \proofstepwidestar[1,2,9,10]{%
    f(\Succ(n))=\eta_R(\Succ(n))}{\rChain{17,21}}
  \proofstepwidestar[1,2,3]{n\in\mathbb N}{\rA}
  \proofstepwidestar[1,2,3]{%
    f(n)=\eta_R(n)}{%
    \FormulaRefAuto{PeanoInductionRule}[ax]{8,22,23}}
\end{tabproofwide}

\FormulaThmDeltaK[Universelle Eigenschaft der natürlichen Zahlen]{%
  \begin{aligned}[t]
    &\SemiringAlg{R}{\oplus}{\odot}{0_R}{1_R}\vdash{}\\[-2pt]
    &\exists! f\,
      \SemiringHom{f}{\mathbb N,+,\cdot,0,1}
        {R,\oplus,\odot,0_R,1_R}
  \end{aligned}
}{NaturalSemiringUniversalHomomorphism}{%
  \DeltaRow{Mengen}{R\dsep0_R\dsep1_R\dsep n}
  \DeltaRow{Funktionen}{f\dsep\oplus\dsep\odot}
}
\begin{tabproofwide}
  \proofstepwidestar[1]{%
    \SemiringAlg{R}{\oplus}{\odot}{0_R}{1_R}}{\rA}
  \proofstepwidestar[1]{%
    \SemiringHom{\NaturalSemiringMap{R,\oplus,\odot,0_R,1_R}}
      {\mathbb N,+,\cdot,0,1}{R,\oplus,\odot,0_R,1_R}}{%
    \FormulaRefAuto{NaturalSemiringCanonicalHomomorphism}[thm]{1}}
  \proofstepwidestar[3]{%
    \SemiringHom{f}{\mathbb N,+,\cdot,0,1}
      {R,\oplus,\odot,0_R,1_R}}{\rA}
  \proofstepwidestar[3]{f\colon\mathbb N\to R}{%
    \FormulaRefAuto{SemiringHomFunctionAxiom}[ax]{3}}
  \proofstepwidestar[1]{%
    \NaturalSemiringMap{R,\oplus,\odot,0_R,1_R}
    \colon\mathbb N\to R}{%
    \FormulaRefAuto{NaturalSemiringMapFunction}[thm]{1}}
  \proofstepwidestar[6]{n\in\mathbb N}{\rA}
  \proofstepwidestar[1,3,6]{%
    f(n)=\NaturalSemiringMap{R,\oplus,\odot,0_R,1_R}(n)}{%
    \FormulaRefAuto{NaturalSemiringHomPointwiseUnique}[thm]{1,3,6}}
  \proofstepwidestar[1,3]{%
    \forall n\in\mathbb N\,
    \bigl(
      f(n)=\NaturalSemiringMap{R,\oplus,\odot,0_R,1_R}(n)
    \bigr)}{%
    \rUI{\rRI{6,7}}}
  \proofstepwidestar[1,3]{%
    f=\NaturalSemiringMap{R,\oplus,\odot,0_R,1_R}}{%
    \FormulaRefAuto{\forall x\in A\,(F(x)=G(x)) \eqvdash F=G}[thm]{8}}
  \proofstepwidestar[1]{%
    \exists! f\,
    \SemiringHom{f}{\mathbb N,+,\cdot,0,1}
      {R,\oplus,\odot,0_R,1_R}}{%
    \UEI{2,3,9}}
\end{tabproofwide}

\section{Surjektivität und Induktion}

\FormulaThmDeltaK[Surjektivität der kanonischen Abbildung genau bei vollem Bild]{%
  \begin{aligned}[t]
    &\SemiringAlg{R}{\oplus}{\odot}{0_R}{1_R}\\
    &\vdash
      \Bigl(
        \NaturalSemiringMap{R,\oplus,\odot,0_R,1_R}
          \colon\mathbb N\sur R
        \leftrightarrow
        \NaturalSemiringMap{R,\oplus,\odot,0_R,1_R}[\mathbb N]=R
      \Bigr)
  \end{aligned}
}{NaturalSemiringMapSurjectiveIffFullImage}{%
  \DeltaRow{Mengen}{R\dsep0_R\dsep1_R}
  \DeltaRow{Funktionen}{\oplus\dsep\odot}
}
\begin{tabproofwide}
  \proofstepwidestar[1]{%
    \SemiringAlg{R}{\oplus}{\odot}{0_R}{1_R}}{\rA}
  \proofstepwidestar[1]{%
    \NaturalSemiringMap{R,\oplus,\odot,0_R,1_R}
      \colon\mathbb N\to R}{%
    \FormulaRefAuto{NaturalSemiringMapFunction}[thm]{1}}
  \proofstepwidestar[1]{%
    \begin{aligned}[t]
      &\NaturalSemiringMap{R,\oplus,\odot,0_R,1_R}
        \colon\mathbb N\sur R\\
      &\leftrightarrow
        \NaturalSemiringMap{R,\oplus,\odot,0_R,1_R}[\mathbb N]=R
    \end{aligned}
  }{%
    \FormulaRefAuto{SurjectiveIffImageEqualsCodomain}[thm]{2}}
\end{tabproofwide}

\FormulaThmDeltaK[Surjektivität der kanonischen Abbildung genau beim Nachfolgerinduktionsprinzip]{%
  \begin{aligned}[t]
    &\SemiringAlg{R}{\oplus}{\odot}{0_R}{1_R}\\
    &\vdash
      \Bigl(
        \NaturalSemiringMap{R,\oplus,\odot,0_R,1_R}
          \colon\mathbb N\sur R
        \leftrightarrow{}\\[-2pt]
    &\qquad
        \forall A\in\powerset(R)\,
        \Bigl(
          \bigl(
            0_R\in A\land
            \forall x\in A\,
            \bigl(
              \SemiringSuccessor{R,\oplus,\odot,0_R,1_R}(x)\in A
            \bigr)
          \bigr)
          \rightarrow A=R
        \Bigr)
      \Bigr)
  \end{aligned}
}{NaturalSemiringMapSurjectiveIffSuccessorInduction}{%
  \DeltaRow{Mengen}{R\dsep A\dsep0_R\dsep1_R\dsep x}
  \DeltaRow{Funktionen}{\oplus\dsep\odot}
}
\begin{notation*}
\[
  \eta_R:=\NaturalSemiringMap{R,\oplus,\odot,0_R,1_R},
  \qquad
  s_R:=\SemiringSuccessor{R,\oplus,\odot,0_R,1_R}.
\]
\end{notation*}
\begin{tabproofwide}
  \proofstepwidestar[1]{%
    \SemiringAlg{R}{\oplus}{\odot}{0_R}{1_R}}{\rA}
  \proofstepwidestar[1]{0_R\in R}{%
    \FormulaRefAuto{SemiringZeroCarrier}[thm]{1}}
  \proofstepwidestar[1]{%
    s_R\colon R\to R}{%
    \FormulaRefAuto{SemiringSuccessorFunction}[thm]{1}}
  \proofstepwidestar[1]{%
    \begin{aligned}[t]
      &\RecFun{0_R}{s_R}\colon\mathbb N\sur R
      \leftrightarrow{}\\[-2pt]
      &\forall A\in\powerset(R)\,
      \Bigl(
        \bigl(
          0_R\in A\land
          \forall x\in A\,
          \bigl(
            s_R(x)\in A
          \bigr)
        \bigr)
        \rightarrow A=R
      \Bigr)
    \end{aligned}
  }{%
    \FormulaRefAuto{RecFunSurjectiveIffInductionPrinciple}[thm]{2,3}}
  \proofstepwidestar[1]{%
    \eta_R=\RecFun{0_R}{s_R}}{%
    \FormulaRefAuto{NaturalSemiringMapDef}[def]{1}}
  \proofstepwidestar[1]{%
    \begin{aligned}[t]
      &\eta_R\colon\mathbb N\sur R
      \leftrightarrow{}\\[-2pt]
      &\forall A\in\powerset(R)\,
      \Bigl(
        \bigl(
          0_R\in A\land
          \forall x\in A\,
          \bigl(
            s_R(x)\in A
          \bigr)
        \bigr)
        \rightarrow A=R
      \Bigr)
    \end{aligned}
  }{\rIE{5,4}}
\end{tabproofwide}

\section{Injektivität}

\FormulaThmDeltaK[Injektivität schließt Kollisionen aus]{%
  \begin{aligned}[t]
    &\SemiringAlg{R}{\oplus}{\odot}{0_R}{1_R}\dsep
      \NaturalSemiringMap{R,\oplus,\odot,0_R,1_R}
        \colon\mathbb N\inj R\\
    &\vdash
      \forall m,n\in\mathbb N\,
      \Bigl(
        \NaturalSemiringMap{R,\oplus,\odot,0_R,1_R}(m)
        =
        \NaturalSemiringMap{R,\oplus,\odot,0_R,1_R}(n)
        \rightarrow m=n
      \Bigr)
  \end{aligned}
}{NaturalSemiringMapInjectiveNoCollision}{%
  \DeltaRow{Mengen}{R\dsep0_R\dsep1_R\dsep m\dsep n}
  \DeltaRow{Funktionen}{\oplus\dsep\odot}
}
\begin{notation*}
\[
  \eta_R:=\NaturalSemiringMap{R,\oplus,\odot,0_R,1_R}.
\]
\end{notation*}
\begin{tabproofwide}
  \proofstepwidestar[1]{%
    \SemiringAlg{R}{\oplus}{\odot}{0_R}{1_R}}{\rA}
  \proofstepwidestar[2]{%
    \eta_R\colon\mathbb N\inj R}{\rA}
  \proofstepwidestar[3]{m\in\mathbb N}{\rA}
  \proofstepwidestar[4]{n\in\mathbb N}{\rA}
  \proofstepwidestar[5]{%
    \eta_R(m)=\eta_R(n)}{\rA}
  \proofstepwidestar[2,3,4,5]{m=n}{%
    \FormulaRefAuto{Injektivität}[ax]{2,3,4,5}}
  \proofstepwidestar[2,3,4]{%
    \begin{aligned}[t]
      &\eta_R(m)=\eta_R(n)\\
      &\rightarrow m=n
    \end{aligned}
  }{\rRI{5,6}}
  \proofstepwidestar[2,3]{%
    \forall n\in\mathbb N\,
    \Bigl(
      \eta_R(m)=\eta_R(n)
      \rightarrow m=n
    \Bigr)}{\rUI{\rRI{4,7}}}
  \proofstepwidestar[2]{%
    \forall m,n\in\mathbb N\,
    \Bigl(
      \eta_R(m)=\eta_R(n)
      \rightarrow m=n
    \Bigr)}{\rUI{\rRI{3,8}}}
\end{tabproofwide}

\FormulaThmDeltaK[Kollisionsfreiheit liefert Injektivität]{%
  \begin{aligned}[t]
    &\SemiringAlg{R}{\oplus}{\odot}{0_R}{1_R}\dsep{}\\[-2pt]
    &\forall m,n\in\mathbb N\,
      \Bigl(
        \NaturalSemiringMap{R,\oplus,\odot,0_R,1_R}(m)
        =
        \NaturalSemiringMap{R,\oplus,\odot,0_R,1_R}(n)
        \rightarrow m=n
      \Bigr)\\
    &\vdash
      \NaturalSemiringMap{R,\oplus,\odot,0_R,1_R}
        \colon\mathbb N\inj R
  \end{aligned}
}{NaturalSemiringMapNoCollisionInjective}{%
  \DeltaRow{Mengen}{R\dsep0_R\dsep1_R\dsep m\dsep n}
  \DeltaRow{Funktionen}{\oplus\dsep\odot}
}
\begin{tabproofwide}
  \proofstepwidestar[1]{%
    \SemiringAlg{R}{\oplus}{\odot}{0_R}{1_R}}{\rA}
  \proofstepwidestar[2]{%
    \forall m,n\in\mathbb N\,
    \Bigl(
      \NaturalSemiringMap{R,\oplus,\odot,0_R,1_R}(m)
      =
      \NaturalSemiringMap{R,\oplus,\odot,0_R,1_R}(n)
      \rightarrow m=n
    \Bigr)}{\rA}
  \proofstepwidestar[1]{%
    \NaturalSemiringMap{R,\oplus,\odot,0_R,1_R}
      \colon\mathbb N\to R}{%
    \FormulaRefAuto{NaturalSemiringMapFunction}[thm]{1}}
  \proofstepwidestar[1,2]{%
    \NaturalSemiringMap{R,\oplus,\odot,0_R,1_R}
      \colon\mathbb N\inj R}{%
    \FormulaRefAuto{Injektive Funktion}[def]{3,2}}
\end{tabproofwide}

\FormulaThmDeltaK[Injektivität der kanonischen Abbildung genau bei Kollisionsfreiheit]{%
  \begin{aligned}[t]
    &\SemiringAlg{R}{\oplus}{\odot}{0_R}{1_R}\\
    &\vdash
      \Bigl(
        \NaturalSemiringMap{R,\oplus,\odot,0_R,1_R}
          \colon\mathbb N\inj R
        \leftrightarrow{}\\[-2pt]
    &\qquad
        \forall m,n\in\mathbb N\,
        \Bigl(
          \NaturalSemiringMap{R,\oplus,\odot,0_R,1_R}(m)
          =
          \NaturalSemiringMap{R,\oplus,\odot,0_R,1_R}(n)
          \rightarrow m=n
        \Bigr)
      \Bigr)
  \end{aligned}
}{NaturalSemiringMapInjectiveIffNoCollision}{%
  \DeltaRow{Mengen}{R\dsep0_R\dsep1_R\dsep m\dsep n}
  \DeltaRow{Funktionen}{\oplus\dsep\odot}
}
\begin{tabproofwide}
    \proofstepwidestar[1]{%
      \SemiringAlg{R}{\oplus}{\odot}{0_R}{1_R}}{\rA}
    \proofstepwidestar[2]{%
      \NaturalSemiringMap{R,\oplus,\odot,0_R,1_R}
        \colon\mathbb N\inj R}{\rA}
    \proofstepwidestar[1,2]{%
      \forall m,n\in\mathbb N\,
      \Bigl(
        \NaturalSemiringMap{R,\oplus,\odot,0_R,1_R}(m)
        =
        \NaturalSemiringMap{R,\oplus,\odot,0_R,1_R}(n)
        \rightarrow m=n
      \Bigr)}{%
      \FormulaRefAuto{NaturalSemiringMapInjectiveNoCollision}[thm]{1,2}}
    \proofstepwidestar[1]{%
      \begin{aligned}[t]
        &\NaturalSemiringMap{R,\oplus,\odot,0_R,1_R}
          \colon\mathbb N\inj R
        \rightarrow{}\\[-2pt]
        &\forall m,n\in\mathbb N\,
        \Bigl(
          \NaturalSemiringMap{R,\oplus,\odot,0_R,1_R}(m)
          =
          \NaturalSemiringMap{R,\oplus,\odot,0_R,1_R}(n)
          \rightarrow m=n
        \Bigr)
      \end{aligned}
    }{\rRI{2,3}}
    \proofstepwidestar[5]{%
      \forall m,n\in\mathbb N\,
      \Bigl(
        \NaturalSemiringMap{R,\oplus,\odot,0_R,1_R}(m)
        =
        \NaturalSemiringMap{R,\oplus,\odot,0_R,1_R}(n)
        \rightarrow m=n
      \Bigr)}{\rA}
    \proofstepwidestar[1,5]{%
      \NaturalSemiringMap{R,\oplus,\odot,0_R,1_R}
        \colon\mathbb N\inj R}{%
      \FormulaRefAuto{NaturalSemiringMapNoCollisionInjective}[thm]{1,5}}
    \proofstepwidestar[1]{%
      \begin{aligned}[t]
        &\forall m,n\in\mathbb N\,
        \Bigl(
          \NaturalSemiringMap{R,\oplus,\odot,0_R,1_R}(m)
          =
          \NaturalSemiringMap{R,\oplus,\odot,0_R,1_R}(n)
          \rightarrow m=n
        \Bigr)\\
        &\rightarrow
        \NaturalSemiringMap{R,\oplus,\odot,0_R,1_R}
          \colon\mathbb N\inj R
      \end{aligned}
    }{\rRI{5,6}}
    \proofstepwidestar[1]{%
      \begin{aligned}[t]
        &\NaturalSemiringMap{R,\oplus,\odot,0_R,1_R}
          \colon\mathbb N\inj R
        \leftrightarrow{}\\[-2pt]
        &\forall m,n\in\mathbb N\,
        \Bigl(
          \NaturalSemiringMap{R,\oplus,\odot,0_R,1_R}(m)
          =
          \NaturalSemiringMap{R,\oplus,\odot,0_R,1_R}(n)
          \rightarrow m=n
        \Bigr)
      \end{aligned}
    }{\rLRI{4,7}}
\end{tabproofwide}

\FormulaThmDeltaK[Peano-Kriterium für die Injektivität der kanonischen Abbildung]{%
  \begin{aligned}[t]
    &\SemiringAlg{R}{\oplus}{\odot}{0_R}{1_R}\dsep
      \SemiringSuccessor{R,\oplus,\odot,0_R,1_R}\colon R\inj R\dsep{}\\[-2pt]
    &\forall x\in R\,
      \Bigl(
        \SemiringSuccessor{R,\oplus,\odot,0_R,1_R}(x)\neq0_R
      \Bigr)\\
    &\vdash
      \NaturalSemiringMap{R,\oplus,\odot,0_R,1_R}
        \colon\mathbb N\inj R
  \end{aligned}
}{NaturalSemiringMapInjectivePeanoCriterion}{%
  \DeltaRow{Mengen}{R\dsep0_R\dsep1_R\dsep x}
  \DeltaRow{Funktionen}{\oplus\dsep\odot}
}
\begin{tabproofwide}
  \proofstepwidestar[1]{%
    \SemiringAlg{R}{\oplus}{\odot}{0_R}{1_R}}{\rA}
  \proofstepwidestar[2]{%
    \SemiringSuccessor{R,\oplus,\odot,0_R,1_R}\colon R\inj R}{\rA}
  \proofstepwidestar[3]{%
    \forall x\in R\,
    \Bigl(
      \SemiringSuccessor{R,\oplus,\odot,0_R,1_R}(x)\neq0_R
    \Bigr)}{\rA}
  \proofstepwidestar[1]{0_R\in R}{%
    \FormulaRefAuto{SemiringZeroCarrier}[thm]{1}}
  \proofstepwidestar[1,2,3]{%
    \RecFun{0_R}
      {\SemiringSuccessor{R,\oplus,\odot,0_R,1_R}}
      \colon\mathbb N\inj R}{%
      \FormulaRefAuto{%
        m_0\in M\dsep f\colon M\inj M\dsep
        \forall x\in M\,\bigl(f(x)\neq m_0\bigr)
        \vdash \RecFun{m_0}{f}\colon\mathbb{N}\inj M}[thm]{4,2,3}}
  \proofstepwidestar[1]{%
    \NaturalSemiringMap{R,\oplus,\odot,0_R,1_R}
    =
    \RecFun{0_R}
      {\SemiringSuccessor{R,\oplus,\odot,0_R,1_R}}}{%
    \FormulaRefAuto{NaturalSemiringMapDef}[def]{1}}
  \proofstepwidestar[1,2,3]{%
    \NaturalSemiringMap{R,\oplus,\odot,0_R,1_R}
      \colon\mathbb N\inj R}{\rIE{6,5}}
\end{tabproofwide}

\FormulaThmDeltaK[Die kanonische Abbildung in einen endlichen Halbring ist nicht injektiv]{%
  \SemiringAlg{R}{\oplus}{\odot}{0_R}{1_R}\dsep\Finite{R}
  \vdash
  \neg\Bigl(
    \NaturalSemiringMap{R,\oplus,\odot,0_R,1_R}
      \colon\mathbb N\inj R
  \Bigr)
}{NaturalSemiringMapNotInjectiveIntoFiniteSemiring}{%
  \DeltaRow{Mengen}{R\dsep0_R\dsep1_R}
  \DeltaRow{Funktionen}{H\dsep\oplus\dsep\odot}
}
\begin{tabproof}
  \proofstep{1}{%
    \SemiringAlg{R}{\oplus}{\odot}{0_R}{1_R}}{\rA}
  \proofstep{2}{\Finite{R}}{\rA}
  \proofstep{2}{%
    \neg\exists H\colon\mathbb N\inj R}{%
    \FormulaRefAuto{FiniteNoNatInjection}[thm]{2}}
  \proofstep{4}{%
    \NaturalSemiringMap{R,\oplus,\odot,0_R,1_R}
      \colon\mathbb N\inj R}{\rA}
  \proofstep{4}{%
    \exists H\colon\mathbb N\inj R}{\rEI{4}}
  \proofstep{2,4}{\bot}{\rBI{3,5}}
  \proofstep{2}{%
    \neg\Bigl(
      \NaturalSemiringMap{R,\oplus,\odot,0_R,1_R}
        \colon\mathbb N\inj R
    \Bigr)}{\rCI{4,6}}
\end{tabproof}

\section{Isomorphie}

\FormulaThmDeltaK[Isomorphie impliziert Injektivität und Surjektivität der kanonischen Abbildung]{%
  \begin{aligned}[t]
    &\SemiringAlg{R}{\oplus}{\odot}{0_R}{1_R}\dsep
      \SemiringIso{\NaturalSemiringMap{R,\oplus,\odot,0_R,1_R}}
        {\mathbb N,+,\cdot,0,1}
        {R,\oplus,\odot,0_R,1_R}\\
    &\vdash
      \NaturalSemiringMap{R,\oplus,\odot,0_R,1_R}
        \colon\mathbb N\inj R
      \land
      \NaturalSemiringMap{R,\oplus,\odot,0_R,1_R}
        \colon\mathbb N\sur R
  \end{aligned}
}{NaturalSemiringMapIsoInjectiveSurjective}{%
  \DeltaRow{Mengen}{R\dsep0_R\dsep1_R}
  \DeltaRow{Funktionen}{\oplus\dsep\odot}
}
\begin{tabproofwide}
    \proofstepwidestar[1]{%
      \SemiringAlg{R}{\oplus}{\odot}{0_R}{1_R}}{\rA}
    \proofstepwidestar[2]{%
      \SemiringIso{\NaturalSemiringMap{R,\oplus,\odot,0_R,1_R}}
        {\mathbb N,+,\cdot,0,1}
        {R,\oplus,\odot,0_R,1_R}}{\rA}
    \proofstepwidestar[2]{%
      \NaturalSemiringMap{R,\oplus,\odot,0_R,1_R}
        \colon\mathbb N\bij R}{%
      \FormulaRefAuto{SemiringIsoDef}[def]{2}}
    \proofstepwidestar[2]{%
      \NaturalSemiringMap{R,\oplus,\odot,0_R,1_R}
        \colon\mathbb N\inj R}{%
      \FormulaRefAuto{%
        F\colon A\bij B\vdash F\colon A\inj B}[thm]{3}}
    \proofstepwidestar[2]{%
      \NaturalSemiringMap{R,\oplus,\odot,0_R,1_R}
        \colon\mathbb N\sur R}{%
      \FormulaRefAuto{%
        F\colon A\bij B\vdash F\colon A\sur B}[thm]{3}}
    \proofstepwidestar[2]{%
      \NaturalSemiringMap{R,\oplus,\odot,0_R,1_R}
        \colon\mathbb N\inj R
      \land
      \NaturalSemiringMap{R,\oplus,\odot,0_R,1_R}
        \colon\mathbb N\sur R}{\rAI{4,5}}
\end{tabproofwide}

\FormulaThmDeltaK[Injektivität und Surjektivität implizieren Isomorphie der kanonischen Abbildung]{%
  \begin{aligned}[t]
    &\SemiringAlg{R}{\oplus}{\odot}{0_R}{1_R}\dsep{}\\[-2pt]
    &\NaturalSemiringMap{R,\oplus,\odot,0_R,1_R}
      \colon\mathbb N\inj R
    \land
    \NaturalSemiringMap{R,\oplus,\odot,0_R,1_R}
      \colon\mathbb N\sur R\\
    &\vdash
      \SemiringIso{\NaturalSemiringMap{R,\oplus,\odot,0_R,1_R}}
        {\mathbb N,+,\cdot,0,1}
        {R,\oplus,\odot,0_R,1_R}
  \end{aligned}
}{NaturalSemiringMapInjectiveSurjectiveIso}{%
  \DeltaRow{Mengen}{R\dsep0_R\dsep1_R}
  \DeltaRow{Funktionen}{\oplus\dsep\odot}
}
\begin{tabproofwide}
    \proofstepwidestar[1]{%
      \SemiringAlg{R}{\oplus}{\odot}{0_R}{1_R}}{\rA}
    \proofstepwidestar[2]{%
      \NaturalSemiringMap{R,\oplus,\odot,0_R,1_R}
        \colon\mathbb N\inj R
      \land
      \NaturalSemiringMap{R,\oplus,\odot,0_R,1_R}
        \colon\mathbb N\sur R}{\rA}
    \proofstepwidestar[2]{%
      \NaturalSemiringMap{R,\oplus,\odot,0_R,1_R}
        \colon\mathbb N\inj R}{\rAEa{2}}
    \proofstepwidestar[2]{%
      \NaturalSemiringMap{R,\oplus,\odot,0_R,1_R}
        \colon\mathbb N\sur R}{\rAEb{2}}
    \proofstepwidestar[2]{%
      \NaturalSemiringMap{R,\oplus,\odot,0_R,1_R}
        \colon\mathbb N\bij R}{%
      \FormulaRefAuto{%
        F\colon A\inj B\dsep F\colon A\sur B
        \vdash F\colon A\bij B}[thm]{3,4}}
    \proofstepwidestar[1]{%
      \SemiringHom{\NaturalSemiringMap{R,\oplus,\odot,0_R,1_R}}
        {\mathbb N,+,\cdot,0,1}
        {R,\oplus,\odot,0_R,1_R}}{%
      \FormulaRefAuto{NaturalSemiringCanonicalHomomorphism}[thm]{1}}
    \proofstepwidestar[1,2]{%
      \SemiringIso{\NaturalSemiringMap{R,\oplus,\odot,0_R,1_R}}
        {\mathbb N,+,\cdot,0,1}
        {R,\oplus,\odot,0_R,1_R}}{%
      \FormulaRefAuto{SemiringIsoDef}[def]{6,5}}
\end{tabproofwide}

\FormulaThmDeltaK[Isomorphie genau bei Injektivität und Surjektivität der kanonischen Abbildung]{%
  \begin{aligned}[t]
    &\SemiringAlg{R}{\oplus}{\odot}{0_R}{1_R}\\
    &\vdash
      \Bigl(
        \SemiringIso{\NaturalSemiringMap{R,\oplus,\odot,0_R,1_R}}
          {\mathbb N,+,\cdot,0,1}
          {R,\oplus,\odot,0_R,1_R}
        \leftrightarrow{}\\[-2pt]
    &\qquad
        \NaturalSemiringMap{R,\oplus,\odot,0_R,1_R}
          \colon\mathbb N\inj R
        \land
        \NaturalSemiringMap{R,\oplus,\odot,0_R,1_R}
          \colon\mathbb N\sur R
      \Bigr)
  \end{aligned}
}{NaturalSemiringMapIsoIffInjectiveAndSurjective}{%
  \DeltaRow{Mengen}{R\dsep0_R\dsep1_R}
  \DeltaRow{Funktionen}{\oplus\dsep\odot}
}
\begin{tabproofwide}
    \proofstepwidestar[1]{%
      \SemiringAlg{R}{\oplus}{\odot}{0_R}{1_R}}{\rA}
    \proofstepwidestar[2]{%
      \SemiringIso{\NaturalSemiringMap{R,\oplus,\odot,0_R,1_R}}
        {\mathbb N,+,\cdot,0,1}
        {R,\oplus,\odot,0_R,1_R}}{\rA}
    \proofstepwidestar[1,2]{%
      \NaturalSemiringMap{R,\oplus,\odot,0_R,1_R}
        \colon\mathbb N\inj R
      \land
      \NaturalSemiringMap{R,\oplus,\odot,0_R,1_R}
        \colon\mathbb N\sur R}{%
      \FormulaRefAuto{NaturalSemiringMapIsoInjectiveSurjective}[thm]{1,2}}
    \proofstepwidestar[1]{%
      \begin{aligned}[t]
        &\SemiringIso{\NaturalSemiringMap{R,\oplus,\odot,0_R,1_R}}
          {\mathbb N,+,\cdot,0,1}
          {R,\oplus,\odot,0_R,1_R}\\
        &\rightarrow
          \NaturalSemiringMap{R,\oplus,\odot,0_R,1_R}
            \colon\mathbb N\inj R
          \land
          \NaturalSemiringMap{R,\oplus,\odot,0_R,1_R}
            \colon\mathbb N\sur R
      \end{aligned}
    }{\rRI{2,3}}
    \proofstepwidestar[5]{%
      \NaturalSemiringMap{R,\oplus,\odot,0_R,1_R}
        \colon\mathbb N\inj R
      \land
      \NaturalSemiringMap{R,\oplus,\odot,0_R,1_R}
        \colon\mathbb N\sur R}{\rA}
    \proofstepwidestar[1,5]{%
      \SemiringIso{\NaturalSemiringMap{R,\oplus,\odot,0_R,1_R}}
        {\mathbb N,+,\cdot,0,1}
        {R,\oplus,\odot,0_R,1_R}}{%
      \FormulaRefAuto{NaturalSemiringMapInjectiveSurjectiveIso}[thm]{1,5}}
    \proofstepwidestar[1]{%
      \begin{aligned}[t]
        &\NaturalSemiringMap{R,\oplus,\odot,0_R,1_R}
          \colon\mathbb N\inj R
        \land
        \NaturalSemiringMap{R,\oplus,\odot,0_R,1_R}
          \colon\mathbb N\sur R\\
        &\rightarrow
          \SemiringIso{\NaturalSemiringMap{R,\oplus,\odot,0_R,1_R}}
            {\mathbb N,+,\cdot,0,1}
            {R,\oplus,\odot,0_R,1_R}
      \end{aligned}
    }{\rRI{5,6}}
    \proofstepwidestar[1]{%
      \begin{aligned}[t]
        &\SemiringIso{\NaturalSemiringMap{R,\oplus,\odot,0_R,1_R}}
          {\mathbb N,+,\cdot,0,1}
          {R,\oplus,\odot,0_R,1_R}
        \leftrightarrow{}\\[-2pt]
        &\NaturalSemiringMap{R,\oplus,\odot,0_R,1_R}
          \colon\mathbb N\inj R
        \land
        \NaturalSemiringMap{R,\oplus,\odot,0_R,1_R}
          \colon\mathbb N\sur R
      \end{aligned}
    }{\rLRI{4,7}}
\end{tabproofwide}

\section{Fehlen additiver Inverser}

\FormulaThmDeltaK[Darstellung als Nachfolger]{%
  n\in\mathbb{N}\vdash1+n=\Succ(n)%
}{NaturalOnePlusAsSuccessor}{%
  \DeltaRow{Mengen}{\mathbb{N}\dsep n}
}
\begin{tabproofwide}
  \proofstepwidestar[1]{n\in\mathbb{N}}{\rA}
  \proofstepwidestar[]{1\in\mathbb{N}}{%
    \FormulaRefAuto{PeanoOneInN}[thm]}
  \proofstepwide[1]{1+n}{=}{n+1}{%
    \FormulaRefAuto{PeanoAddCommutative}[thm]{2,1}}
  \proofstepwide[1]{}{=}{\Succ(n)}{%
    \FormulaRefAuto{PeanoAddOneSucc}[thm]{1}}
  \proofstepwidestar[1]{1+n=\Succ(n)}{%
    \rChain{3,4}}
\end{tabproofwide}

\FormulaThmDeltaK[Widerspruch zum Nullaxiom]{%
  \exists n\in\mathbb{N}\,(1+n=0)\vdash\bot%
}{NaturalAdditionNotGroupContradiction}{%
  \DeltaRow{Mengen}{\mathbb{N}\dsep n}
}
\begin{tabproofwide}
  \proofstepwidestar[1]{%
    \exists n\in\mathbb{N}\,(1+n=0)}{\rA}
  \proofstepwidestar[2]{n\in\mathbb{N}}{\rA}
  \proofstepwidestar[3]{1+n=0}{\rA}
  \proofstepwidestar[2]{1+n=\Succ(n)}{%
    \FormulaRefAuto{NaturalOnePlusAsSuccessor}{2}}
  \proofstepwide[2]{\Succ(n)}{=}{1+n}{%
    \FormulaRefAuto{a=b\vdash b=a}{4}}
  \proofstepwide[2,3]{}{=}{0}{\rIE{3}}
  \proofstepwidestar[2,3]{\Succ(n)=0}{\rChain{5,6}}
  \proofstepwidestar[2]{\Succ(n)\neq0}{%
    \FormulaRefAuto{PeanoSuccNonzero}[ax]{2}}
  \proofstepwidestar[2,3]{\bot}{\rCE{8,7}}
  \proofstepwidestar[1]{\bot}{\rEE{1,2,3,9}}
\end{tabproofwide}

\FormulaThmDeltaK[Die Eins besitzt kein additives Inverses in \(\mathbb{N}\)]{%
  \neg\exists n\in\mathbb{N}\,(1+n=0)%
}{NaturalAdditionNotGroup}{%
  \DeltaRow{Mengen}{\mathbb{N}\dsep n}
}
\begin{tabproof}
  \proofstep{1}{\exists n\in\mathbb{N}\,(1+n=0)}{\rA}
  \proofstep{1}{\bot}{%
    \FormulaRefAuto{NaturalAdditionNotGroupContradiction}[thm]{1}}
  \proofstep{}{\neg\exists n\in\mathbb{N}\,(1+n=0)}{\rCI{1,2}}
\end{tabproof}

\FormulaThmDeltaK[Spezialisierung auf die Eins]{%
  \begin{aligned}[t]
    &\forall a\in\mathbb{N}\,
      \exists b\in\mathbb{N}\,(a+b=0\land b+a=0)\\
    &\vdash\exists b\in\mathbb{N}\,(1+b=0)
  \end{aligned}
}{NaturalRingInverseRequirementAtOne}{%
  \DeltaRow{Mengen}{\mathbb{N}\dsep a\dsep b}
}
\begin{tabproofwide}
  \proofstepwidestar[1]{%
    \forall a\in\mathbb{N}\,
    \exists b\in\mathbb{N}\,(a+b=0\land b+a=0)}{\rA}
  \proofstepwidestar[]{1\in\mathbb{N}}{%
    \FormulaRefAuto{PeanoOneInN}[thm]}
  \proofstepwidestar[1]{%
    \exists b\in\mathbb{N}\,(1+b=0\land b+1=0)}{%
    \rRE{2,\rUE{1}}}
  \proofstepwidestar[4]{b\in\mathbb{N}}{\rA}
  \proofstepwidestar[5]{1+b=0\land b+1=0}{\rA}
  \proofstepwidestar[5]{1+b=0}{\rAEa{5}}
  \proofstepwidestar[4,5]{%
    \exists b\in\mathbb{N}\,(1+b=0)}{\rEI{4,6}}
  \proofstepwidestar[1]{%
    \exists b\in\mathbb{N}\,(1+b=0)}{\rEE{3,4,5,7}}
\end{tabproofwide}

\FormulaThmDeltaK[Fehlschlagen der Inversenforderung]{%
  \neg\forall a\in\mathbb{N}\,
  \exists b\in\mathbb{N}\,(a+b=0\land b+a=0)%
}{NaturalNumbersNotRingConclusion}{%
  \DeltaRow{Mengen}{\mathbb{N}\dsep a\dsep b}
}
\begin{tabproofwide}
  \proofstepwidestar[]{%
    \neg\exists b\in\mathbb{N}\,(1+b=0)}{%
    \FormulaRefAuto{NaturalAdditionNotGroup}[thm]}
  \proofstepwidestar[2]{%
    \forall a\in\mathbb{N}\,
    \exists b\in\mathbb{N}\,(a+b=0\land b+a=0)}{\rA}
  \proofstepwidestar[2]{%
    \exists b\in\mathbb{N}\,(1+b=0)}{%
    \FormulaRefAuto{NaturalRingInverseRequirementAtOne}[thm]{2}}
  \proofstepwidestar[2]{\bot}{\rCE{1,3}}
  \proofstepwidestar[]{%
    \neg\forall a\in\mathbb{N}\,
    \exists b\in\mathbb{N}\,(a+b=0\land b+a=0)}{%
    \rCI{2,4}}
\end{tabproofwide}

\FormulaThmDeltaK[Die natürlichen Zahlen erfüllen nicht die additive Inversenforderung eines Rings]{%
  \begin{aligned}[t]
    &\CommutativeSemiringAlg{\mathbb{N}}{+}{\cdot}{0}{1}\land{}\\[-2pt]
    &\neg\forall a\in\mathbb{N}\,
      \exists b\in\mathbb{N}\,(a+b=0\land b+a=0)
  \end{aligned}
}{NaturalNumbersNotRing}{%
  \DeltaRow{Mengen}{\mathbb{N}\dsep a\dsep b}
}
\begin{tabproofwide}
  \proofstepwidestar[]{%
    \CommutativeSemiringAlg{\mathbb{N}}{+}{\cdot}{0}{1}}{%
    \FormulaRefAuto{NaturalCommutativeUnitalSemiring}[thm]}
  \proofstepwidestar[]{%
    \neg\forall a\in\mathbb{N}\,
    \exists b\in\mathbb{N}\,(a+b=0\land b+a=0)%
  }{\FormulaRefAuto{NaturalNumbersNotRingConclusion}[thm]}
  \proofstepwidestar[]{%
    \CommutativeSemiringAlg{\mathbb{N}}{+}{\cdot}{0}{1}\land
    \neg\forall a\in\mathbb{N}\,
    \exists b\in\mathbb{N}\,(a+b=0\land b+a=0)}{%
    \rAI{1,2}}
\end{tabproofwide}

\end{document}
